
\documentclass{article}

% Recommended, but optional, packages for figures and better typesetting:
\usepackage{microtype}
\usepackage{graphicx}
\usepackage{subcaption}
\usepackage{booktabs} % for professional tables

% hyperref makes hyperlinks in the resulting PDF.
% If your build breaks (sometimes temporarily if a hyperlink spans a page)
% please comment out the following usepackage line and replace
% \usepackage{icml2026} with \usepackage[nohyperref]{icml2026} above.
% \usepackage{hyperref}



% Attempt to make hyperref and algorithmic work together better:
\newcommand{\theHalgorithm}{\arabic{algorithm}}

% Use the following line for the initial blind version submitted for review:
\usepackage{icml2026}

% For preprint, use
% \usepackage[preprint]{icml2026}

% If accepted, instead use the following line for the camera-ready submission:
% \usepackage[accepted]{icml2026}


% Additional packages (avoiding duplicates with template packages above)
\usepackage[utf8]{inputenc} % allow utf-8 input
\usepackage[T1]{fontenc}% use 8-bit T1 fonts
\usepackage{url}% simple URL typesetting
\usepackage{amsfonts}% blackboard math symbols
\usepackage{nicefrac}% compact symbols for 1/2, etc.
\usepackage{xcolor}% colors

\usepackage{multirow}
\usepackage{rotating}
\usepackage{siunitx}
\usepackage{makecell}
\usepackage{bm}
\usepackage{amsmath}
\usepackage{amssymb}
\usepackage{mathtools}
\usepackage{amsthm}
\usepackage{gensymb}
% \usepackage{algorithm}
% \usepackage[noend]{algpseudocode}
\newcommand{\return}{\textbf{return }}
\newcommand{\Comment}[1]{\hfill $\triangleright$ #1}
\usepackage{enumitem}

\usepackage{xcolor}
\definecolor{cvprblue}{rgb}{0.21,0.49,0.74}
\usepackage[pagebackref,breaklinks,colorlinks,citecolor=cvprblue]{hyperref}
% Add a period to the end of an abbreviation unless there's one
% already, then \xspace.
\usepackage{xspace}
\makeatletter
\DeclareRobustCommand\onedot{\futurelet\@let@token\@onedot}
\def\@onedot{\ifx\@let@token.\else.\null\fi\xspace}

\def\eg{\emph{e.g}\onedot} \def\Eg{\emph{E.g}\onedot}
\def\ie{\emph{i.e}\onedot} \def\Ie{\emph{I.e}\onedot}
\def\cf{\emph{c.f}\onedot} \def\Cf{\emph{C.f}\onedot}
\def\etc{\emph{etc}\onedot} \def\vs{\emph{vs}\onedot}
\def\wrt{w.r.t\onedot} \def\dof{d.o.f\onedot}
\def\aka{\emph{a.k.a}\onedot}
\def\etal{\emph{et al}\onedot}
\makeatother


\definecolor{DarkMagenta}{rgb}{0.7, 0.0, 0.7}
\newcommand{\mznote}[1]{\textcolor{green}{\textbf{MZ: #1}}}
\newcommand{\nguyen}[1]{{\color{DarkMagenta}#1}}
\newcommand{\nguyenrmk}[1]{{\color{DarkMagenta} {\bf [VN: #1]}}}
\newcommand{\zgnote}[1]{{\color{blue} {\bf [ZG: #1]}}}

\newcommand{\beginsupplement}{%
\setcounter{table}{0}
\renewcommand{\thetable}{S\arabic{table}}%
\setcounter{figure}{0}
\renewcommand{\thefigure}{S\arabic{figure}}%
}

% Custom commands
\newcommand{\modelname}{\textit{d}GS}

% Note: \onedot, abbreviation commands (\eg, \ie, \etc, \etal, \wrt, \dof, \aka)
% are already defined earlier in this file
% Note: amsmath, amssymb, mathtools, amsthm already loaded above (lines 44-47)

% if you use cleveref..
\usepackage[capitalize,noabbrev]{cleveref}

%%%%%%%%%%%%%%%%%%%%%%%%%%%%%%%%
% THEOREMS
%%%%%%%%%%%%%%%%%%%%%%%%%%%%%%%%
\theoremstyle{plain}
\newtheorem{theorem}{Theorem}[section]
\newtheorem{proposition}[theorem]{Proposition}
\newtheorem{lemma}[theorem]{Lemma}
\newtheorem{corollary}[theorem]{Corollary}
\theoremstyle{definition}
\newtheorem{definition}[theorem]{Definition}
\newtheorem{assumption}[theorem]{Assumption}
\theoremstyle{remark}
\newtheorem{remark}[theorem]{Remark}

% Todonotes is useful during development; simply uncomment the next line
%    and comment out the line below the next line to turn off comments
%\usepackage[disable,textsize=tiny]{todonotes}
\usepackage[textsize=tiny]{todonotes}

% The \icmltitle you define below is probably too long as a header.
% Therefore, a short form for the running title is supplied here:
\icmltitlerunning{\modelname: Direct Parameterization for Efficient Conditional Gaussian Splatting}

\begin{document}

\twocolumn[
  \icmltitle{\modelname: Direct Parameterization for Efficient Conditional Gaussian Splatting}

  % It is OKAY to include author information, even for blind submissions: the
  % style file will automatically remove it for you unless you've provided
  % the [accepted] option to the icml2026 package.

  % List of affiliations: The first argument should be a (short) identifier you
  % will use later to specify author affiliations Academic affiliations
  % should list Department, University, City, Region, Country Industry
  % affiliations should list Company, City, Region, Country

  % You can specify symbols, otherwise they are numbered in order. Ideally, you
  % should not use this facility. Affiliations will be numbered in order of
  % appearance and this is the preferred way.
  \icmlsetsymbol{equal}{*}

  \begin{icmlauthorlist}
    \icmlauthor{Zhongpai Gao}{uii}
    % Additional co-authors can be added here if needed:
    % \icmlauthor{Meng Zheng}{uii}
    % \icmlauthor{Benjamin Planche}{uii}
    % \icmlauthor{Anwesa Choudhuri}{uii}
    % \icmlauthor{Terrence Chen}{uii}
    % \icmlauthor{Ziyan Wu}{uii}
  \end{icmlauthorlist}

  \icmlaffiliation{uii}{United Imaging Intelligence, Boston, MA, USA}

  \icmlcorrespondingauthor{Zhongpai Gao}{zhongpai.gao@uii-ai.com}

  % You may provide any keywords that you find helpful for describing your
  % paper; these are used to populate the "keywords" metadata in the PDF but
  % will not be shown in the document
  \icmlkeywords{Gaussian Splatting, Neural Rendering, Novel View Synthesis, 3D Reconstruction}

  \vskip 0.3in
{%
\begin{center}
  \centering
    \captionsetup{type=figure}
   \includegraphics[width=0.95\textwidth]{dGS/figures/teaser.pdf}
    \captionof{figure}{\textbf{Rendering quality and efficiency comparison.} \modelname{} achieves better quality and higher efficiency than 6DGS \citep{gao20246dgs} on the static \texttt{bicycle} scene (left) and 7DGS \citep{gao20257dgs} on the dynamic \texttt{heart1} scene (right).}
   \label{fig:teaser}
\end{center}
}
]


% this must go after the closing bracket ] following \twocolumn[ ...

% This command actually creates the footnote in the first column listing the
% affiliations and the copyright notice. The command takes one argument, which
% is text to display at the start of the footnote. The \icmlEqualContribution
% command is standard text for equal contribution. Remove it (just {}) if you
% do not need this facility.

% Use ONE of the following lines. DO NOT remove the command.
% If you have no special notice, KEEP empty braces:
\printAffiliationsAndNotice{}  % no special notice (required even if empty)
% Or, if applicable, use the standard equal contribution text:
% \printAffiliationsAndNotice{\icmlEqualContribution}
\begin{abstract}
   High-dimensional Gaussian Splatting (N-DGS) methods like 6DGS and 7DGS achieve impressive view- and temporal-dependent modeling through conditional N-dimensional Gaussians. However, their monolithic covariance design derives all conditioning effects from a joint covariance matrix, requiring matrix inversions, regression multiplications, and covariance corrections per Gaussian in every forward pass. We introduce \modelname{}, which replaces this implicit derivation with direct parameterization: opacity modulation via Cholesky precision factors and position shifts via displacement matrices with learnable per-dimension coupling strengths. This eliminates matrix inversions and covariance corrections, achieving 5--6$\times$ faster slicing kernels and up to 2.66$\times$ faster end-to-end rendering. Across five benchmarks spanning static, dynamic, and specular scenes, \modelname{} consistently matches or exceeds N-DGS quality on every dataset (up to +0.97 dB on Mip-NeRF 360), offering a principled, efficient alternative to covariance-derived conditioning in high-dimensional Gaussian splatting.
\end{abstract}

\section{Introduction}
\label{sec:introduction}
3D Gaussian Splatting (3DGS)~\citep{3dgs} revolutionized novel-view synthesis by achieving photorealistic quality at real-time speeds through explicit 3D Gaussian primitives. However, real-world scenes exhibit complex phenomena beyond static representations: specular reflections, temporal dynamics, and view-dependent effects that require richer modeling capacity.

Recent extensions elevate Gaussians into higher-dimensional spaces to capture these phenomena. N-dimensional Gaussians (N-DG)~\citep{diolatzis2024n} pioneered this direction, refined by 6D Gaussian Splatting (6DGS)~\citep{gao20246dgs} which models position and viewing direction jointly for view-dependent effects, and 7D Gaussian Splatting (7DGS)~\citep{gao20257dgs} which adds temporal dimensions for dynamic scenes (collectively referred to as N-DGS). Alternative kernel designs like Deformable Beta Splatting (DBS)~\citep{liu2025deformable} improve geometric fidelity but remain limited to spatial dimensions.


Despite their expressive power, these N-dimensional formulations suffer from a fundamental limitation: their covariance-based parameterization. All conditioning effects (position shifts, opacity modulation, and covariance adjustments) are derived from a single joint covariance matrix via conditional Gaussian formulas. Computing any single effect requires inverting the conditioning block, performing regression matrix multiplications, and computing covariance corrections per Gaussian, creating computational overhead that accumulates across millions of primitives. Moreover, this implicit derivation offers limited control: position and opacity conditioning are rigidly tied through the shared covariance structure, preventing independent adjustment of their coupling strengths. This is problematic because different phenomena require different conditioning strategies. For instance, view-dependent reflections primarily need opacity modulation, while temporal motion requires position displacement, yet N-DGS's fixed coupling forces the same relationship for all Gaussians. The monolithic design also risks numerical instabilities when covariance matrices become ill-conditioned.

\textbf{Our Insight.} Conditional Gaussian formulas are elegant mathematically but computationally inefficient. Each conditioning effect can be \textbf{directly parameterized} with efficient, independent parameters instead of being implicitly derived through the covariance matrix. Opacity modulation only needs a precision matrix for distance computation; therefore, we parameterize it directly via Cholesky factors. Position shifts only need a displacement matrix scaled by query deviations; therefore, we parameterize it directly with bounded weights. By replacing implicit derivation with explicit parameters, we eliminate matrix inversions while gaining fine-grained control over each effect independently.

We introduce \textbf{\modelname{}}, a direct parameterization for efficient conditional Gaussian splatting. The `\textbf{d}' signifies our principle: \textbf{d}irect parameterization---explicit \textbf{d}isplacement matrices for position and \textbf{d}ecay factors for opacity, enabling \textbf{d}ynamic Gaussian behavior. Like N-DGS, we enable Gaussians to adaptively shift position and modulate opacity based on queries, but replace implicit covariance derivation with efficient explicit parameters. Our key contributions are:


\begin{itemize}[leftmargin=*]
\item \textbf{Direct Opacity Parameterization:} We directly parameterize opacity decay via Cholesky precision matrices, eliminating the need for matrix inversions while guaranteeing numerical stability through triangular operations.

\item \textbf{Direct Position Displacement:} Position shifts are directly parameterized via displacement matrices with spatial scaling and bounded activations, ensuring tight displacement bounds without requiring covariance inversion.

\item \textbf{Learnable Coupling Strength:} Per-dimension coupling parameters control how strongly position shifts respond to query deviations, enabling adaptive interpolation from no position conditioning to full coupling through the shared precision matrix.

\item \textbf{Efficient CUDA Implementation:} We implement optimized CUDA kernels for both our method and N-DGS baseline, fusing conditional operations with register-level computation for fair comparison (Note that the original N-DGS only implements in PyTorch).
\end{itemize}

Experiments demonstrate that \modelname{} matches or exceeds N-DGS quality across all benchmarks while achieving up to 2.66$\times$ faster rendering (166\% throughput increase) on in-the-wild scenes (Mip-NeRF 360).


\section{Related Work}
\label{sec:related_work}

\textbf{3D Gaussian Splatting.}
3D Gaussian Splatting (3DGS)~\citep{3dgs} achieves real-time novel view synthesis through explicit 3D Gaussian primitives with spherical harmonic color encoding. This has inspired extensions addressing multi-scale anti-aliasing~\citep{yu2023mipsplattingaliasfree3dgaussian,yan2024multiscale3dgaussiansplatting}, compression~\cite{lee2024compact,morgenstern2024compact}, and semantic modeling~\cite{shi2023languageembedded3dgaussians,zhou2024feature3dgssupercharging3d}. However, 3DGS remains limited to static geometry, relying on spherical harmonics for view-dependent appearance without mechanisms for temporal dynamics or view-dependent geometry.

\textbf{Alternative Kernel Designs.}
Several works explore non-Gaussian kernels to overcome limitations in 3D spatial modeling. Modified Gaussian variants include generalized exponentials~\cite{hamdi2024gesgeneralizedexponentialsplatting}, truncated kernels~\cite{liu2025tnt}, second-order formulations~\cite{zhang2025quadratic}, half-Gaussians~\cite{li20253d}, and discontinuity-aware profiles~\cite{qu2024disc}. Alternative probability distributions have also been explored: Beta kernels~\cite{liu2025deformable} for improved geometric fidelity, Student-t distributions~\cite{zhu20253dstudentsplattingscooping} for heavier tails, and Gabor functions~\cite{zhou20253dgabsplat3dgaborsplatting} for sharper features. Beyond volumetric representations, methods using convex shapes~\cite{held20243dconvexsplattingradiance}, triangles~\cite{held2025triangle}, and radial kernels~\cite{huang2025deformable} offer alternative geometric primitives. However, all these approaches focus on 3D spatial kernels and do not extend to higher dimensions for modeling view-dependent or temporal effects.

\textbf{Higher-Dimensional Gaussian Splatting.}
To model view-dependent and temporal effects, recent methods lift Gaussians into higher-dimensional spaces. 4DGS~\cite{yang2023real} extends 3DGS with temporal dimensions for dynamic scene reconstruction. N-DG~\cite{diolatzis2024n} establishes a general framework for arbitrary-dimensional Gaussians. Building on this, 6DGS~\cite{gao20246dgs} jointly models spatial position and viewing direction for view-dependent rendering, while 7DGS~\cite{gao20257dgs} unifies spatial, angular, and temporal dimensions. We refer to 6DGS and 7DGS collectively as N-DGS. UBS~\citep{liu2025universal} extends this paradigm by applying Beta kernels to N-dimensional representations. A key limitation is that N-DGS methods derive all conditioning effects from joint covariance matrices, requiring matrix inversions, regression multiplications, and covariance corrections per primitive, adding significant per-Gaussian overhead and risking numerical instability.

\textbf{Dynamic Scene Modeling.}
NeRF-based methods model dynamics through deformations: D-NeRF~\cite{pumarola2020dnerfneuralradiancefields} introduced canonical-to-deformed mappings, Nerfies~\cite{park2021nerfies} handled non-rigid motion, and HyperNeRF~\cite{park2021hypernerf} captured topological changes. For Gaussian Splatting, 4DGS~\cite{yang2023real} enables real-time dynamic rendering, while Dynamic 3D Gaussians~\cite{luiten2023dynamic} combines canonical representations with learned deformations. FreeTimeGS~\cite{wang2025freetimegs} allows Gaussians to appear at arbitrary times and locations with motion functions. These deformation-based approaches model dynamics as transformations of primitives, whereas 7DGS~\cite{gao20257dgs} achieves unified modeling through higher-dimensional conditioning.


\section{Preliminary}

% Our work builds upon prior Gaussian Splatting methods. To contextualize our contribution, we first detail the progression from static 3D models to the more complex 6D and 7D formulations, highlighting the trade-offs between expressive power and computational cost at each step.

\paragraph{3D Gaussian Splatting}
3D Gaussian Splatting (3DGS)~\cite{3dgs} revolutionized novel view synthesis by representing scenes as a collection of anisotropic 3D Gaussians. Each Gaussian is a primitive defined by a set of learnable attributes that describe its geometry and appearance. Its geometry is determined by a mean vector \(\boldsymbol{\mu} \in \mathbb{R}^3\) specifying its center in 3D space, and a covariance matrix \(\boldsymbol{\Sigma} \in \mathbb{R}^{3 \times 3}\) defining its shape, scale, and orientation. \(\boldsymbol{\Sigma}\) is parameterized by a diagonal scaling matrix \(\mathbf{S} = \operatorname{diag}(s_x, s_y, s_z)\) and a rotation matrix \(\mathbf{R}\) (represented as a quaternion), such that \(\boldsymbol{\Sigma} = \mathbf{R} \mathbf{S} \mathbf{S}^\top \mathbf{R}^\top\). The appearance of each Gaussian is governed by a scalar opacity \(\alpha \in [0, 1]\) that controls its transparency, and a view-dependent color modeled using Spherical Harmonics (SH). The color \(c\) for a given viewing direction \(d\) is computed as a sum over the SH basis functions \(Y_{\ell m}\):
\begin{equation}
c(d) = \sum_{\ell=0}^{L} \sum_{m=-\ell}^{\ell} \beta_{\ell m} Y_{\ell m}(d),
\end{equation}
where \(\beta_{\ell m}\) are the learnable SH coefficients. 3DGS achieves state-of-the-art quality and real-time rendering for static scenes. However, its representation is fundamentally static and cannot model geometric changes arising from complex view-dependent effects or temporal dynamics.

\paragraph{N-Dimensional Gaussian Splatting}

To capture view-dependent and temporal effects, recent methods extend 3DGS into higher dimensions: 6DGS~\cite{gao20246dgs} for view-dependent phenomena and 7DGS~\cite{gao20257dgs} for dynamic scenes. Each primitive is an N-dimensional Gaussian over position $\mathbf{X}_p \in \mathbb{R}^3$ and query variables $\mathbf{X}_q \in \mathbb{R}^C$ (view direction for 6DGS with $C=3$, view+time for 7DGS with $C=4$):
\begin{align}
\mathbf{X} = \begin{pmatrix} \mathbf{X}_p \\ \mathbf{X}_q \end{pmatrix} \sim \mathcal{N}\!\left(
\begin{pmatrix}
\boldsymbol{\mu}_p \\ \boldsymbol{\mu}_q
\end{pmatrix},
\begin{pmatrix}
\boldsymbol{\Sigma}_{pp} & \boldsymbol{\Sigma}_{pq} \\
\boldsymbol{\Sigma}_{pq}^\top & \boldsymbol{\Sigma}_{qq}
\end{pmatrix}
\right).
\end{align}
At query $\mathbf{q}$, conditioning yields a 3D Gaussian with parameters:
\begin{align}
\alpha_{\text{cond}} &= \alpha \cdot \exp\left(-\lambda\,(\mathbf{q}-\boldsymbol{\mu}_q)^\top \boldsymbol{\Sigma}_{qq}^{-1}(\mathbf{q}-\boldsymbol{\mu}_q)\right),\\
\boldsymbol{\mu}_{\text{cond}} &= \boldsymbol{\mu}_p + \boldsymbol{\Sigma}_{pq}\,\boldsymbol{\Sigma}_{qq}^{-1}(\mathbf{q}-\boldsymbol{\mu}_q), \\
\boldsymbol{\Sigma}_{\text{cond}} &= \boldsymbol{\Sigma}_{pp} - \boldsymbol{\Sigma}_{pq}\,\boldsymbol{\Sigma}_{qq}^{-1}\,\boldsymbol{\Sigma}_{pq}^\top,
\end{align}
where $\lambda$ is a scaling factor (default 0.35). This formulation requires inverting $\boldsymbol{\Sigma}_{qq}$, computing the regression matrix $\boldsymbol{\Sigma}_{pq}\boldsymbol{\Sigma}_{qq}^{-1}$, and evaluating the covariance correction $\boldsymbol{\Sigma}_{pq}\boldsymbol{\Sigma}_{qq}^{-1}\boldsymbol{\Sigma}_{pq}^\top$ for every Gaussian in every forward pass. While $C$ is small (3 or 4), these three matrix operations collectively add significant per-Gaussian overhead that accumulates across millions of primitives and risks numerical instability during training, motivating our simpler direct parameterization.



\begin{figure*}[t]
\centering

\includegraphics[width=0.99\linewidth]{figures/pipeline.pdf} % Ensure this figure exists
\caption{Overview of the \modelname{} pipeline. Given query deviation $\boldsymbol{\delta} = q - \mu_q$ from the canonical view/time, we directly parameterize two conditioning effects: (1) \textbf{Opacity modulation} via Cholesky factor $\mathbf{L}$, and (2) \textbf{Position displacement} via displacement matrix $\mathbf{V}_{pq}$ and scaled precision $\tilde{\mathbf{V}}_{qq}$ with learnable coupling $\boldsymbol{\Lambda}$. The spatial covariance $\boldsymbol{\Sigma}_{pp}$ remains query-independent. This eliminates matrix inversions and covariance corrections while enabling learnable control over position coupling strength via $\boldsymbol{\Lambda}$.}
\label{fig:architecture}
\vspace{-1em}
\end{figure*}

\section{Our Approach}

N-DGS methods derive all conditioning effects from a single $(3+C) \times (3+C)$ joint covariance matrix via conditional Gaussian formulas. To compute any effect (position shift, opacity modulation, or covariance adjustment), they must invert $\boldsymbol{\Sigma}_{qq}$, multiply regression matrices, and compute covariance corrections, requiring multiple matrix operations per Gaussian per forward pass.

We propose \modelname{}, which replaces this \textbf{implicit derivation} with \textbf{direct parameterization}. Rather than computing conditioning effects through matrix algebra, we explicitly parameterize each component with dedicated, efficient parameters. This direct approach eliminates matrix inversions and covariance corrections while providing independent control over each conditioning effect.

\subsection{Overview}

Our approach directly parameterizes conditional Gaussian slicing with two explicit components. Given the deviation $\boldsymbol{\delta} = q - \mu_q$ from the canonical view/time:
\begin{align}
\alpha_{\text{cond}} &= \alpha \cdot \exp\left(-\|\mathbf{L}^\top \boldsymbol{\delta}\|_2^2\right), & &\text{(Opacity)} \\
\mu_{\text{cond}} &= \mu_p + \mathbf{V}_{pq} \cdot \tilde{\mathbf{V}}_{qq} \cdot \boldsymbol{\delta}, & &\text{(Position)} \end{align}
where $\mathbf{L} \in \mathbb{R}^{C \times C}$ is the directly parameterized Cholesky factor, $\mathbf{V}_{qq} = \mathbf{L}\mathbf{L}^\top$ is the precision matrix computed directly from $\mathbf{L}$ without inversion, $\mathbf{V}_{pq} = \bar{s}\boldsymbol{\Theta}_{pq} \in \mathbb{R}^{3 \times C}$ is the displacement matrix (mean spatial scale $\bar{s}$ times normalized $\boldsymbol{\Theta}_{pq}$ with $\|\boldsymbol{\Theta}_{pq}\|_F = 1$), and $\tilde{\mathbf{V}}_{qq} = \text{diag}(\boldsymbol{\Lambda}) \mathbf{V}_{qq}$ is the scaled precision with learnable per-dimension coupling strengths $\boldsymbol{\Lambda} \in [0,1]^C$ (dimensions of the same type share parameters). The spatial covariance $\boldsymbol{\Sigma}_{pp}$ remains independent of the query, using the standard 3DGS quaternion-scale parameterization ($\boldsymbol{\Sigma}_{pp} = \mathbf{R}\mathbf{S}\mathbf{S}^\top \mathbf{R}^\top$) for seamless compatibility with existing pipelines, as densification strategies rely on the spatial covariance.

For dynamic scenes with view and time conditioning ($C=4$), we further \textbf{separate view and time} by using a block-diagonal precision matrix with zero cross-terms:
\begin{equation}
\mathbf{V}_{qq} = \mathbf{L}\mathbf{L}^\top = \begin{pmatrix}
\mathbf{V}_{vv} & \mathbf{0} \\
\mathbf{0} & v_t
\end{pmatrix},
\end{equation}
where $\mathbf{V}_{vv} \in \mathbb{R}^{3 \times 3}$ is the view precision and $v_t \in \mathbb{R}$ is the time precision. This separation eliminates coupling between spatial view directions and temporal dynamics, reflecting the physical intuition that view-dependent effects and temporal changes are independent phenomena.


\begin{table}[t!]
\centering
\caption{Per-Gaussian operations: N-DGS (implicit) vs \modelname{} (direct). In N-DGS, $\mathbf{K} = \boldsymbol{\Sigma}_{pq}\boldsymbol{\Sigma}_{qq}^{-1}$ is reused for position shift and covariance correction. \modelname{} eliminates inversion, regression, and covariance correction.}
\label{tab:comparison_overview}
\setlength{\tabcolsep}{4pt}
\resizebox{0.48\textwidth}{!}{%
\begin{tabular}{@{}l|ccc@{}}
\toprule
\textbf{Operation} & \textbf{N-DGS} & \textbf{\textit{d}GS-O} & \textbf{\textit{d}GS} \\
\midrule
Invert $\boldsymbol{\Sigma}_{qq}$ ($C{\times}C$) & $\boldsymbol{\Sigma}_{qq}^{-1}$ & --- & --- \\
Regr.\ matrix $\mathbf{K}$ ($3{\times}C \cdot C{\times}C$) & $\boldsymbol{\Sigma}_{pq}\boldsymbol{\Sigma}_{qq}^{-1}$ & --- & --- \\
Cov.\ correction ($3{\times}C \cdot C{\times}3$) & $\mathbf{K}\,\boldsymbol{\Sigma}_{pq}^\top$ & --- & --- \\
Position shift & $\mathbf{K}\,\boldsymbol{\delta}$ & --- & $\mathbf{V}_{pq} \tilde{\mathbf{V}}_{qq} \boldsymbol{\delta}$ \\
Opacity & $\boldsymbol{\delta}^\top \boldsymbol{\Sigma}_{qq}^{-1} \boldsymbol{\delta}$ & $\|\mathbf{L}^\top \boldsymbol{\delta}\|^2$ & $\|\mathbf{L}^\top \boldsymbol{\delta}\|^2$ \\
\bottomrule
\end{tabular}%
}
\vspace{-1em}
\end{table}

Table~\ref{tab:comparison_overview} compares the per-Gaussian operations of N-DGS and \modelname{}. N-DGS must compute the regression matrix $\mathbf{K} = \boldsymbol{\Sigma}_{pq}\boldsymbol{\Sigma}_{qq}^{-1}$, which is reused for position shift and covariance correction; \modelname{} eliminates all three operations. The key innovation is \textbf{learnable coupling strength} via per-dimension parameters $\Lambda$: when $\Lambda = \mathbf{0}$, position becomes query-independent ($\boldsymbol{\mu}_{\text{cond}} = \boldsymbol{\mu}_p$), which we denote as the simplified \textbf{\textit{d}GS-O} variant using only direct opacity parameterization for maximum efficiency; when $\Lambda = \mathbf{1}$, position couples through the shared precision matrix $\mathbf{V}_{qq}$ like N-DGS. The following subsections detail the direct parameterization of each component.

\subsection{Opacity Conditioning}

In N-DGS, opacity is modulated as $\alpha_{\text{cond}} = \alpha \cdot \exp(-\lambda\,(\mathbf{q}-\boldsymbol{\mu}_q)^\top \boldsymbol{\Sigma}_{qq}^{-1}(\mathbf{q}-\boldsymbol{\mu}_q))$, requiring matrix inversion.

Our formulation directly parameterizes the precision matrix via its Cholesky factor $\mathbf{L}$, computing $\mathbf{L}^\top \boldsymbol{\delta}$ followed by $\|\cdot\|_2^2$ through a triangular matrix-vector multiplication. This eliminates the need for matrix inversion. Parameterizing the Cholesky factor directly with exponential activation on the diagonal guarantees that $\mathbf{V}_{qq} = \mathbf{L}\mathbf{L}^\top$ is always symmetric positive definite, avoiding numerical instabilities from ill-conditioned matrices.

For dynamic scenes, the block-diagonal structure enables \textbf{separated view and time opacity}. With $\boldsymbol{\delta} = (\boldsymbol{\delta}_v, \delta_t)^\top$ where $\boldsymbol{\delta}_v \in \mathbb{R}^3$ is the view deviation and $\delta_t \in \mathbb{R}$ is the time deviation, the opacity factors as:
\begin{equation}
\alpha_{\text{cond}} = \alpha \cdot \exp\left(-\left(\boldsymbol{\delta}_v^\top \mathbf{V}_{vv} \boldsymbol{\delta}_v + \delta_t^2 v_t\right)\right).
\end{equation}
This factorization allows view- and time-dependent opacity to be controlled independently, with view precision $\mathbf{V}_{vv}$ parameterized by a $3 \times 3$ Cholesky factor (6 parameters) and time precision $v_t$ by a single scalar, for a total of 7 parameters instead of 10 for a full $4 \times 4$ precision matrix.

\subsection{Position Conditioning}
\label{subsec:position}

In N-DGS, position is computed as $\boldsymbol{\mu}_{\text{cond}} = \boldsymbol{\mu}_p + \boldsymbol{\Sigma}_{pq}\boldsymbol{\Sigma}_{qq}^{-1}(\mathbf{q}-\boldsymbol{\mu}_q)$, where both $\boldsymbol{\Sigma}_{pq}$ and $\boldsymbol{\Sigma}_{qq}^{-1}$ are implicitly derived from the joint covariance matrix through Cholesky decomposition and matrix inversion.

We replace this implicit derivation with \textbf{direct parameterization} of position shifts, combined with learnable coupling strength control:
\begin{equation}
\vspace{-0.3em}
\boldsymbol{\mu}_{\text{cond}} = \boldsymbol{\mu}_p + \underbrace{\bar{s} \boldsymbol{\Theta}_{pq}}_{\mathbf{V}_{pq}} \cdot \underbrace{\text{diag}(\Lambda) \mathbf{V}_{qq}}_{\tilde{\mathbf{V}}_{qq}} \cdot \boldsymbol{\delta},
\vspace{-0.3em}
\end{equation}
where $\bar{s} = \text{mean}(s_x, s_y, s_z)$ is the spatial scale, $\boldsymbol{\Theta}_{pq} \in \mathbb{R}^{3 \times C}$ are learnable parameters normalized to unit Frobenius norm ($\|\boldsymbol{\Theta}_{pq}\|_F = 1$), $\text{diag}(\boldsymbol{\Lambda}) \in \mathbb{R}^{C \times C}$ contains learnable per-dimension coupling strengths in $[0,1]$, and $\mathbf{V}_{qq} = \mathbf{L}\mathbf{L}^\top$ is the precision matrix shared with opacity conditioning. The coupling strength $\boldsymbol{\Lambda}$ controls how strongly position shifts respond to query deviations:
\vspace{-0.5em}
\begin{itemize}[leftmargin=*]
\item When $\Lambda = \mathbf{0}$: $\tilde{\mathbf{V}}_{qq} = \mathbf{0}$, position becomes query-independent ($\boldsymbol{\mu}_{\text{cond}} = \boldsymbol{\mu}_p$, no conditioning).
\item When $\Lambda = \mathbf{1}$: $\tilde{\mathbf{V}}_{qq} = \mathbf{V}_{qq}$, position and opacity couple through the shared precision matrix $\mathbf{V}_{qq}$ (similar to N-DGS, but with direct $\mathbf{V}_{pq}$ instead of derived $\boldsymbol{\Sigma}_{pq}$).
\item Intermediate values: Partial coupling with per-dimension control. Different query dimensions (e.g., view vs. time) can have different coupling strengths, with dimensions of the same type (e.g., the three view dimensions) sharing parameters.
\vspace{-0.5em}
\end{itemize}
This direct parameterization eliminates covariance inversion while providing fine-grained control over position conditioning strength. Each Gaussian learns optimal per-dimension coupling during training.

Inspired by N-DGS's design principle that displacement should scale with spatial extent (so small Gaussians produce small displacements), our displacement matrix $\mathbf{V}_{pq} = \bar{s} \boldsymbol{\Theta}_{pq}$ goes beyond their implicit covariance-based approach through direct parameterization. By normalizing $\boldsymbol{\Theta}_{pq}$ to unit Frobenius norm ($\|\boldsymbol{\Theta}_{pq}\|_F = 1$), displacements scale proportionally with the average Gaussian extent $\bar{s} = \text{mean}(s_x, s_y, s_z)$, with each row uniformly bounded by $\|\mathbf{V}_{pq}^{(i)}\|_2 \leq \bar{s}$. In contrast, N-DGS's implicit Cholesky structure yields a looser displacement bound of $\sqrt{C}(i{-}1 + s_i)$ per spatial row $i$, where the constant offset $(i{-}1)\sqrt{C}$ can dominate for small Gaussians and introduces asymmetry across spatial dimensions depending on dimension ordering. Our direct parameterization provides a tighter bound ($\bar{s}$) with no offset and uniform treatment across all spatial dimensions, better handling anisotropic Gaussians. (See Appendix~\ref{sec:spatial_scaling_proof} for the full derivation.)


\begin{algorithm}[tb]
\caption{\textit{d}GS Conditional Gaussian Slicing}
\label{alg:slice}
\begin{algorithmic}[1]
\REQUIRE Gaussian parameters: $\mu_p$, $\mu_q$, $\alpha$, $\mathbf{S}$, $\boldsymbol{\Theta}_{pq}$, $\mathbf{L}$, $\Lambda$
\REQUIRE Query point: $q$ (view/time)
\STATE $\boldsymbol{\delta} \gets q - \mu_q$ %\Comment{Compute query deviation}
\STATE $\alpha_{\text{cond}} \gets \alpha \cdot \exp(-\|\mathbf{L}^\top \boldsymbol{\delta}\|^2)$ % \Comment{Conditional opacity}
\IF{\textit{d}GS-O variant}
    \STATE \return  $\alpha_{\text{cond}}, \mu_p$ \Comment{Early exit for opacity only}
\ENDIF
\STATE $\mathbf{V}_{pq} \gets \text{mean}(\text{diag}(\mathbf{S})) \cdot \boldsymbol{\Theta}_{pq}$ % \Comment{Construct displacement matrix}
\STATE $\tilde{\mathbf{V}}_{qq} \gets \text{diag}(\boldsymbol{\Lambda}) \mathbf{L}\mathbf{L}^\top$ % \Comment{Scaled precision for coupling}
\STATE $\mu_{\text{cond}} \gets \mu_p + \mathbf{V}_{pq} \cdot \tilde{\mathbf{V}}_{qq} \cdot \boldsymbol{\delta}$ % \Comment{Conditional position}
\STATE \return $\alpha_{\text{cond}}$, $\mu_{\text{cond}}$
\end{algorithmic}
\end{algorithm}
\setlength{\textfloatsep}{22pt}


\subsection{CUDA Implementation}

Since the original N-DGS methods, \ie, 6DGS \cite{gao20246dgs} and 7DGS \cite{gao20257dgs}, use PyTorch's autograd for conditional Gaussian operations without custom CUDA kernels, we implement both our disentangled formulation and the N-DGS baseline with optimized CUDA kernels for fair comparison. Our implementation uses a custom kernel that fuses conditional position and opacity computations in a single pass, as shown in Algorithm~\ref{alg:slice}.

\textbf{CUDA Kernel Design.} The kernel is templated on dimension $C$ (3 for 6D static scenes, 4 for 7D dynamic scenes) with compile-time loop unrolling for optimal performance. Each thread processes one Gaussian primitive independently, enabling perfect parallelization across millions of Gaussians. For opacity computation (Line 2 in Algorithm~\ref{alg:slice}), we compute the triangular matrix-vector product $\mathbf{L}^\top \boldsymbol{\delta}$ directly in registers using the lower-triangular structure—only $C(C+1)/2$ multiplications are needed, followed by a dot product for $\|\cdot\|_2^2$. For position displacement (Lines 6-8), we construct $\mathbf{V}_{pq}$ via diagonal scaling of $\tanh$ activations, then compute the scaled precision matrix $\tilde{\mathbf{V}}_{qq}$ by applying per-dimension scaling $\text{diag}(\Lambda)$ to $\mathbf{L}\mathbf{L}^\top$. All matrix operations use registers rather than shared memory, as the small matrix sizes fit entirely in per-thread registers. This fused design minimizes global memory bandwidth by avoiding intermediate tensor allocations and maximizes GPU occupancy through register-level computation.

\begin{figure*}[!t]
    \centering
\includegraphics[width=0.85\linewidth]{dGS/figures/qualitative.pdf}
    \caption{Qualitative comparison of methods for static and dynamic scenes (zoom in for details).}
    \label{fig:qualitative}
    \vspace{-1em}
\end{figure*}

\textbf{Modular Compatibility.} Our implementation provides compile-time options to enable different levels of conditioning: (1) standard 3DGS with no conditioning, (2) \textit{d}GS-O with opacity modulation only, and (3) full \textit{d}GS with both opacity and position modulation. This modularity maintains compatibility with existing 3DGS codebases while allowing practitioners to select the appropriate complexity-performance tradeoff for their application.

\textbf{Kernel Benchmarks.} As shown in Table~\ref{tab:comparison_overview}, N-DGS requires inversion, regression, and covariance correction per Gaussian, while \modelname{} eliminates all three. While each individual operation is small for $C{=}3$ or $4$, eliminating them collectively removes the dominant cost of the slicing kernel. To isolate the slicing speedup from other pipeline components (rasterization, densification), we benchmark the raw CUDA kernels (C++ forward pass, no autograd overhead) on an A100 GPU. Table~\ref{tab:kernel_benchmark} shows that our slicing kernel is consistently \textbf{5--6$\times$ faster} (\textit{d}GS) and \textbf{8--11$\times$ faster} (\textit{d}GS-O) than N-DGS across both $C{=}3$ and $C{=}4$ configurations and scales from 500K to 2M Gaussians. The consistent speedup confirms this is a structural advantage from eliminating matrix inversions and covariance corrections, not a measurement artifact. End-to-end rendering achieves up to 2.66$\times$ speedup (Table~\ref{tab:summary_results}), with the remaining time dominated by rasterization shared across all methods.

\begin{table}[t]
\centering
\caption{CUDA kernel benchmark: slicing-only forward pass time (ms) on A100. Both N-DGS and \modelname{} use the same C++ entry point with identical inputs, isolating the slicing computation.}
\label{tab:kernel_benchmark}
\setlength{\tabcolsep}{4pt}
\resizebox{0.48\textwidth}{!}{%
\begin{tabular}{@{}c|r|rrr|cc@{}}
\toprule
& \#Gauss & N-DGS & \textit{d}GS-O & \textit{d}GS & N-DGS/\textit{d}GS-O & N-DGS/\textit{d}GS \\
\midrule
\multirow{3}{*}{$C{=}3$}
& 500K & 0.44 & 0.05 & 0.09 & \textbf{9.5$\times$} & \textbf{5.1$\times$} \\
& 1M & 0.85 & 0.11 & 0.16 & \textbf{8.0$\times$} & \textbf{5.2$\times$} \\
& 2M & 1.63 & 0.20 & 0.31 & \textbf{8.1$\times$} & \textbf{5.3$\times$} \\
\midrule
\multirow{2}{*}{$C{=}4$}
& 500K & 0.77 & 0.07 & 0.16 & \textbf{11.1$\times$} & \textbf{4.8$\times$} \\
& 1M & 1.40 & 0.15 & 0.23 & \textbf{9.4$\times$} & \textbf{6.2$\times$} \\
\bottomrule
\end{tabular}}
\vspace{-1em}
\end{table}

\begin{table*}[t]
\centering
\caption{Comparison of methods on three static datasets with MCMC densification. Mean values across all scenes are reported (\#Gauss in K, training time in minutes). FPS and FPS$_{\text{CUDA}}$ represent that conditional Gaussian operations are implemented in PyTorch and CUDA, respectively. Note that, N-DGS, \textit{d}GS-O, and \textit{d}GS are integrated with TC-GS \citep{liao2025tc} to speed up rendering.}
\label{tab:mcmc_comparison}
\resizebox{0.72\textwidth}{!}{%
\begin{tabular}{@{}cl|rrrr|ccc@{}}
\toprule
Dataset & Method & \#Gauss & Time $\downarrow$ & FPS $\uparrow$& FPS$_{\text{CUDA}}$ $\uparrow$& PSNR $\uparrow$ & SSIM $\uparrow$ & LPIPS $\downarrow$ \\
\midrule
\multirow{4}{*}{\makecell{\texttt{NeRF}\\\texttt{Synthetic}}}
& 3DGS & 300 & \bf 7.80 & 303.09 & --- & 33.72 & 0.971 & 0.028 \\
& N-DGS & 300 & 10.43 & 179.12 & 348.95 & 34.04 & \bf 0.973 & 0.027 \\
& \textit{d}GS-O (\textbf{Ours}) & 300 & 9.71 & \bf 400.96 & 587.84 & 34.25 & \bf 0.973 & \bf 0.026 \\
& \textit{d}GS (\textbf{Ours}) & 300 & 10.26 & 267.33 & \bf 593.94 & \bf 34.34 & \bf 0.973 & \bf 0.026 \\
\midrule
\multirow{4}{*}{\makecell{\texttt{Mip-NeRF}\\\texttt{360}}}
& 3DGS & 3,104 & \bf 40.16 & 47.62 & --- & 27.06 & 0.788 & 0.237 \\
& N-DGS & 2,852 & 51.23 & 33.79 & 75.31 & 26.88 & 0.775 & 0.263 \\
& \textit{d}GS-O (\textbf{Ours}) & 2,852 & 48.11 & \bf 91.43 & 200.17 & 27.36 & 0.788 & 0.242 \\
& \textit{d}GS (\textbf{Ours}) & 2,574 & 49.85 & 43.72 & \bf 200.56 & \bf 27.67 & \bf 0.794 & \bf 0.227 \\
\midrule
\multirow{4}{*}{\texttt{6DGS-PBR}}
& 3DGS & 163 & 53.28& 183.78 & --- & 30.89 & 0.938 & 0.077 \\
& N-DGS & 163 & 54.80 & 292.60 & 476.79 & 40.06 & \bf 0.976 & 0.039 \\
& \textit{d}GS-O (\textbf{Ours}) & 163 & 48.72  & \bf 467.06 & \bf 588.23& 39.51 & 0.971 & 0.046 \\
& \textit{d}GS (\textbf{Ours}) & 163 & \bf 47.54 & 364.74 & 551.02 & \bf 40.34 & \bf 0.976 & \bf 0.038 \\
\bottomrule
\end{tabular}}
\end{table*}


\begin{table*}[h]
\centering
\caption{Comparison with N-DGS on three static and two dynamic datasets with standard 3DGS densification. Mean values across all scenes are reported (\#Gauss in K, training time in minutes). For comparisons with 3DGS \cite{3dgs} and 4DGS \cite{yang2023real} baselines, refer to the original N-DGS~\cite{gao20246dgs,gao20257dgs}. See Appendix~\ref{sec:detailed_results} for per-scene details.}
\label{tab:summary_results}
\resizebox{0.72\textwidth}{!}{%
\begin{tabular}{@{}ccl|rrr|ccc@{}}
\toprule
& Dataset & Method & \#Gauss & Time $\downarrow$ & FPS$_{\text{CUDA}}$ $\uparrow$ & PSNR $\uparrow$ & SSIM $\uparrow$ & LPIPS $\downarrow$ \\
\midrule
\multirow{9}{*}{\rotatebox{90}{STATIC}}
& \multirow{3}{*}{\makecell{\texttt{NeRF}\\\texttt{Synthetic}}}
& N-DGS & \bf 171 & 6.35 & 615.57 & 33.38 & 0.970 & 0.032 \\
&& \textit{d}GS-O (\textbf{Ours}) & 201 & \bf 5.83 & \bf 754.98 & 33.84 & 0.971 & 0.030 \\
&& \textit{d}GS (\textbf{Ours}) & 201 & 6.84 & 704.84 & \bf 33.91 & \bf 0.972 & \bf 0.029 \\
\cmidrule{2-9}
& \multirow{3}{*}{\makecell{\texttt{Mip-NeRF}\\\texttt{360}}}
& N-DGS & \bf 1273 & 27.09 & 147.46 & 26.40 & 0.739 & 0.316 \\
&& \textit{d}GS-O (\textbf{Ours}) & 2363 & \bf 25.58 & \bf 279.30 & 27.18 & 0.779 & 0.252 \\
&& \textit{d}GS (\textbf{Ours}) & 2463 & 28.21 & 237.82 & \bf 27.37 & \bf 0.781 & \bf 0.251 \\
\cmidrule{2-9}
& \multirow{3}{*}{\texttt{6DGS-PBR}}
& N-DGS & \bf 84 & \bf 26.77 & 678.40 & 37.11 & 0.971 & 0.054 \\
&& \textit{d}GS-O (\textbf{Ours}) & 118 & 30.72 & \bf 696.31 & 36.54 & 0.965 & 0.059 \\
&& \textit{d}GS (\textbf{Ours}) & 134 & 33.80 & 653.44 & \bf 37.75 & \bf 0.972 & \bf 0.052 \\
\midrule
\midrule
\multirow{6}{*}{\rotatebox{90}{DYNAMIC}}
& \multirow{3}{*}{\texttt{D-NeRF}}
& N-DGS & \bf 69 & 14.08 & 852.01 & 33.13 & 0.968 & 0.038 \\
&& \textit{d}GS-O (\textbf{Ours}) & 70 & \bf 13.20 & \bf 948.77 & 30.92 & 0.959 & 0.044 \\
&& \textit{d}GS (\textbf{Ours}) & 85 & 15.56 & 818.41 & \bf 33.59 & \bf 0.971 & \bf 0.032 \\
\cmidrule{2-9}
& \multirow{3}{*}{\texttt{7DGS-PBR}}
& N-DGS & \bf 162 & 38.14 & 709.83 & 31.75 & 0.955 & 0.055 \\
&& \textit{d}GS-O (\textbf{Ours}) & 227 & \bf 34.58 & \bf 835.24 & 31.70 & 0.955 & 0.053 \\
&& \textit{d}GS (\textbf{Ours}) & 226 & 36.96 & 753.77 & \bf 32.18 & \bf 0.956 & \bf 0.050 \\
\bottomrule
\end{tabular}}
\vspace{-0.5em}
\end{table*}

\section{Experiments}

\subsection{Experimental Setup}

\textbf{Datasets.}
We evaluate \modelname{} across five diverse datasets covering both static and dynamic scenes:
\vspace{-1em}
\begin{itemize}[leftmargin=*]
    \item \textit{Static scenes:} (1) \textit{NeRF Synthetic}~\cite{mildenhall2021nerf} with 8 synthetic scenes under controlled lighting for evaluating geometric reconstruction; (2) \textit{Mip-NeRF 360}~\cite{barron2021mipnerf} with 9 challenging unbounded real-world scenes for evaluating view synthesis; and (3) \textit{6DGS-PBR}~\cite{gao20246dgs} featuring 7 physically-based scenes with complex view-dependent effects.

    \item \textit{Dynamic scenes:} (4) \textit{D-NeRF}~\cite{pumarola2020dnerfneuralradiancefields} with 8 synthetic sequences featuring various deformation types; and (5) \textit{7DGS-PBR}~\cite{gao20257dgs} containing 6 physically-based dynamic scenes with complex spatio-temporal-angular correlations.
\end{itemize}


\textbf{Evaluation Metrics.}
We assess reconstruction quality using Peak Signal-to-Noise Ratio (PSNR), Structural Similarity Index (SSIM)~\cite{wang2004image}, and Learned Perceptual Image Patch Similarity (LPIPS)~\cite{zhang2018perceptual}. We also report rendering speed (FPS), training time (minutes), and Gaussian count (K) to evaluate computational efficiency.

\textbf{Implementation Details.}
All methods are implemented in PyTorch on a single NVIDIA A100 GPU. We compare three configurations: 1) \textbf{N-DGS}: full N-DGS formulation with joint position and opacity conditioning, considered as our entangled baseline, 2) \textbf{\textit{d}GS-O}: our disentangled approach with only opacity conditioning ($\mu_{\text{cond}} = \mu_p$, no position displacement), and 3) \textbf{\textit{d}GS}: our full disentangled approach with both opacity conditioning and learnable position coupling via $\Lambda$. We use the same training settings with Adam optimizer for 30K iterations for all three configurations. For initialization, we set the view mean $\mu_d$ as random unit vectors (with time uniformly sampled in $[0, 1]$ for dynamic scenes), initialize displacement parameters $\boldsymbol{\Theta}_{pq}$ from $\mathcal{N}(0, 0.01^2)$, and set the Cholesky factor $\mathbf{L}$ with diagonal elements $\log(2.0)$ and off-diagonal elements 0.

\subsection{Results with MCMC Densification}

Table~\ref{tab:mcmc_comparison} compares our \modelname{} variants against N-DGS \cite{gao20246dgs} and 3DGS \cite{3dgs} baselines on static scene datasets with MCMC densification~\cite{kheradmand20243d}. Following the original paper settings, we set the maximum number of Gaussians to 300K for NeRF Synthetic scenes, and use scene-specific caps for Mip-NeRF 360 (ranging from 1.5M to 6M) and 6DGS-PBR (ranging from 60K to 300K) depending on scene complexity. Notably, our approach demonstrates significantly larger efficiency gains on large-scale scenes with millions of Gaussians. On Mip-NeRF 360 (1.5M-6M Gaussians), both \textit{d}GS-O and \textit{d}GS achieve substantial +166\% FPS speedups over N-DGS (200.17 vs 75.31 FPS and 200.56 vs 75.31 FPS, respectively), while improving quality (+0.48 dB and +0.79 dB PSNR). In contrast, on the smaller-scale 6DGS-PBR dataset (60K-300K Gaussians), speedups are more modest: +23\% for \textit{d}GS-O (588.23 vs 476.79 FPS) and +16\% for \textit{d}GS (551.02 vs 476.79 FPS), with \textit{d}GS improving quality by +0.28 dB. On NeRF Synthetic (300K Gaussians), both \textit{d}GS-O and \textit{d}GS outperform N-DGS in quality (+0.21 and +0.30 dB PSNR) while achieving significant speedups (+68\% and +70\% FPS).

\textbf{CUDA Implementation Benefits.} Comparing FPS vs FPS$_{\text{CUDA}}$ in Table~\ref{tab:mcmc_comparison} demonstrates the substantial performance gains from our optimized CUDA kernels. For N-DGS, CUDA implementation achieves 1.95$\times$ speedup on NeRF Synthetic (348.95 vs 179.12 FPS), 2.23$\times$ on Mip-NeRF 360 (75.31 vs 33.79 FPS), and 1.63$\times$ on 6DGS-PBR (476.79 vs 292.60 FPS). Our \textit{d}GS method shows even more dramatic improvements: 2.22$\times$ speedup on NeRF Synthetic (593.94 vs 267.33 FPS), 4.59$\times$ on Mip-NeRF 360 (200.56 vs 43.72 FPS), and 1.51$\times$ on 6DGS-PBR (551.02 vs 364.74 FPS). The larger speedups on Mip-NeRF 360 highlight how CUDA optimization becomes increasingly critical for large-scale scenes with millions of Gaussians, where our direct parameterization's elimination of matrix inversions maximally benefits from GPU parallelization.

\subsection{Results with Standard Densification}

Table \ref{tab:summary_results} summarizes quality metrics (PSNR, SSIM, LPIPS) and efficiency metrics (Gaussian count, training time, FPS) across all five datasets with standard 3DGS densification \citep{3dgs}, including both static and dynamic scenes. Per-scene results are provided in Appendix~\ref{sec:detailed_results}.

\textbf{Static Scenes.}
On NeRF Synthetic, both variants improve quality: \textit{d}GS-O by +0.46 dB and \textit{d}GS by +0.53 dB, with rendering speedups of 23\% and 15\% respectively. On Mip-NeRF 360, efficiency gains are most pronounced: \textit{d}GS-O achieves 89\% faster rendering with +0.78 dB improvement, while \textit{d}GS achieves 61\% faster rendering with +0.97 dB improvement. Notably, \textit{d}GS-O uses 86\% more primitives (2363K vs 1273K) yet remains faster, demonstrating that per-primitive computation cost dominates over primitive count. On 6DGS-PBR, \textit{d}GS-O shows modest 3\% speedup but -0.57 dB quality drop, while \textit{d}GS recovers and exceeds baseline quality (+0.64 dB).

\textbf{Dynamic Scenes.}
Position conditioning proves essential for temporal dynamics. On D-NeRF, \textit{d}GS-O suffers -2.21 dB quality drop despite 11\% faster rendering, while \textit{d}GS with learnable coupling via $\Lambda$ exceeds N-DGS by +0.46 dB, demonstrating adaptive coupling surpasses fixed coupling. On 7DGS-PBR combining dynamics with angular lighting, \textit{d}GS-O achieves 18\% faster rendering with minimal quality loss (-0.05 dB), while \textit{d}GS improves quality by +0.43 dB with 6\% faster rendering.

\textbf{Gaussian Count.} Under standard densification, \modelname{} retains more primitives than N-DGS (e.g., 2463K vs 1273K on Mip-NeRF 360). This reflects N-DGS's tendency to over-prune on challenging scenes: on \texttt{stump}, N-DGS prunes to 644K / 23.32 dB while \textit{d}GS retains 3343K / 26.00 dB. Importantly, the MCMC results in Table~\ref{tab:mcmc_comparison} cap all methods at the same Gaussian budget, and \textit{d}GS still outperforms N-DGS by +0.79 dB on Mip-NeRF 360, confirming that quality gains are from the parameterization, not additional primitives.

Figure~\ref{fig:qualitative} presents qualitative comparisons. Our \textit{d}GS improves novel view synthesis quality with fewer artifacts.

\begin{table}[!t]
\centering
\caption{Ablation study on NeRF Synthetic and 6DGS-PBR datasets. $\Lambda{=}0$ corresponds to \textit{d}GS-O (opacity only), $\Lambda{=}1$ fixes full coupling, and learned $\Lambda$ is the full \textit{d}GS. ``w/o $\bar{s}$'' removes spatial scaling from displacement.}
\label{tab:ablation}
\setlength{\tabcolsep}{4pt}
\resizebox{0.48\textwidth}{!}{%
\begin{tabular}{l|ccc|ccc}
\toprule
 & \multicolumn{3}{c|}{\texttt{NeRF Synthetic}} & \multicolumn{3}{c}{\texttt{6DGS-PBR}} \\
Config & PSNR $\uparrow$ & SSIM $\uparrow$ & LPIPS $\downarrow$ & PSNR $\uparrow$ & SSIM $\uparrow$ & LPIPS $\downarrow$ \\
\midrule
$\Lambda{=}0$ (\textit{d}GS-O) & 33.84 & 0.971 & 0.030 & 36.54 & 0.965 & 0.059 \\
$\Lambda{=}1$ (fixed) & 33.78 & 0.971 & 0.029 & 37.31 & \textbf{0.972} & 0.053 \\
\textit{d}GS w/o $\bar{s}$ & 31.26 & 0.957 & 0.046 & 34.78 & 0.951 & 0.075 \\
$\Lambda$ learned (\textit{d}GS) & \textbf{33.91} & \textbf{0.972} & \textbf{0.029} & \textbf{37.75} & \textbf{0.972} & \textbf{0.052} \\
\bottomrule
\end{tabular}}%
\vspace{-1em}
\end{table}

\subsection{Ablation Study}

To understand the contribution of each component, we conduct an ablation study comparing four configurations in Table~\ref{tab:ablation}: (1) \textbf{$\Lambda{=}0$} (\textit{d}GS-O): opacity conditioning only, (2) \textbf{$\Lambda{=}1$}: fixed full coupling, (3) \textbf{\textit{d}GS w/o $\bar{s}$}: learned $\Lambda$ without spatial scaling, and (4) \textbf{learned $\Lambda$} (\textit{d}GS): our full method.

\textbf{Learned $\Lambda$ outperforms both fixed extremes.} On NeRF Synthetic, learned $\Lambda$ achieves 33.91 dB vs 33.84 ($\Lambda{=}0$) and 33.78 ($\Lambda{=}1$). On 6DGS-PBR (view-dependent scenes), the gap widens: 37.75 dB vs 36.54 ($\Lambda{=}0$, +1.21 dB) and 37.31 ($\Lambda{=}1$, +0.44 dB). This confirms that optimal coupling is scene-dependent: on diffuse scenes $\Lambda{=}0$ slightly outperforms $\Lambda{=}1$, while on specular scenes $\Lambda{=}1$ is substantially better (+0.77 dB), validating that per-Gaussian adaptive coupling is the right design.

\textbf{Spatial scaling is critical.} Removing $\bar{s}$ from displacement degrades PSNR by 2.65 dB on NeRF Synthetic (33.91 vs 31.26) and 2.97 dB on 6DGS-PBR (37.75 vs 34.78). Position displacement without proper scaling (\textit{d}GS w/o $\bar{s}$) performs worse than no displacement at all ($\Lambda{=}0$), confirming that coupling displacement magnitude to Gaussian spatial extent is essential.

\section{Limitations and Future Work}

\textbf{Limitations.} Under standard densification, \modelname{} tends to retain more Gaussian primitives than N-DGS (e.g., 2463K vs 1273K on Mip-NeRF 360). While these additional primitives contribute positively to reconstruction quality and rendering remains faster due to per-primitive compute savings, they increase memory consumption. Under MCMC densification with equal Gaussian budgets, \modelname{} still outperforms N-DGS (e.g., +0.79 dB on Mip-NeRF 360), confirming that quality gains are not solely from additional primitives. Additionally, \modelname{} is currently evaluated only within the N-DGS Gaussian framework. Combining our direct parameterization with alternative kernels (e.g., Beta kernels in UBS~\citep{liu2025universal}) is a promising direction that could bring efficiency gains to a broader class of N-dimensional splatting methods. We present a concrete formulation for this combination (dBS: direct Beta Splatting) in Appendix~\ref{sec:dbs}, along with a direct comparison between \modelname{} and UBS on shared benchmarks.

\textbf{Future Work.} Our direct parameterization framework naturally extends to additional conditioning effects beyond opacity and position. Analogous to factorized attention mechanisms where Gaussian attributes (opacity, position) serve as \textbf{values} independently modulated by contextual variables (view direction, time) as \textbf{queries}, our approach provides systematic extensibility along two axes: the \textit{value space} (which Gaussian attributes can be conditioned, such as \textbf{rotation} or \textbf{color}) and the \textit{query space} (which variables drive conditioning, such as \textbf{lighting direction} or \textbf{material properties}). Each new conditioning effect can be directly parameterized with dedicated parameters and learnable coupling strength, avoiding the growing overhead in N-DGS where adding dimensions increases the cost of matrix inversions, regression multiplications, and covariance corrections. This modular design ensures per-Gaussian cost scales linearly with conditioning effects.

\section{Conclusion}

We presented \modelname{}, a direct parameterization for efficient N-dimensional Gaussian splatting. Our key insight is that conditional Gaussian formulas, while elegant, are computationally inefficient: each conditioning effect can be directly parameterized with dedicated parameters instead of being implicitly derived through covariance matrix. This direct approach delivers efficiency by eliminating matrix inversions and covariance corrections (5-6$\times$ faster slicing kernel), stability through explicit Cholesky parameterization, and flexibility via learnable per-dimension coupling strengths that control how strongly each effect responds to queries. Experiments demonstrate up to 2.66$\times$ faster rendering on in-the-wild scenes (Mip-NeRF 360) while matching or exceeding N-DGS quality across diverse benchmarks.

% References section
\bibliographystyle{icml2026}
\bibliography{references}

%%%%%%%%%%%%%%%%%%%%%%%%%%%%%%%%%%%%%%%%%%%%%%%%%%%%%%%%%%%%%%%%%%%%%%%%%%%%%%%
%%%%%%%%%%%%%%%%%%%%%%%%%%%%%%%%%%%%%%%%%%%%%%%%%%%%%%%%%%%%%%%%%%%%%%%%%%%%%%%
% APPENDIX
%%%%%%%%%%%%%%%%%%%%%%%%%%%%%%%%%%%%%%%%%%%%%%%%%%%%%%%%%%%%%%%%%%%%%%%%%%%%%%%
%%%%%%%%%%%%%%%%%%%%%%%%%%%%%%%%%%%%%%%%%%%%%%%%%%%%%%%%%%%%%%%%%%%%%%%%%%%%%%%
\newpage
\appendix
\onecolumn
\section{Comparison to N-DGS Displacement Parameterization}
\label{sec:spatial_scaling_proof}

Both N-DGS \cite{gao20246dgs, gao20257dgs} and \modelname{} scale displacement with spatial scale $s_i$, but differ in how tightly this coupling is achieved and in their parameterization approach.

\textbf{N-DGS: Implicit coupling via Cholesky structure.} N-DGS parameterizes the full $(3+C) \times (3+C)$ covariance via Cholesky decomposition $\boldsymbol{\Sigma} = \mathbf{L}\mathbf{L}^\top$, where $\mathbf{L}$ is lower triangular with diagonal elements $L_{ii} = s_i$ (spatial and query scales) and off-diagonal elements $L_{ij} = \ell_{ij}$ for $i > j$, where $\ell_{ij} \in [-1,1]$ are bounded correlation parameters (initialized to zero).

To illustrate this structure, consider 6DGS with 3 spatial and 3 view dimensions ($C=3$). The Cholesky factor $\mathbf{L} \in \mathbb{R}^{6 \times 6}$ has the form:
\begin{equation}
\mathbf{L} = \left(\begin{array}{rrrrrr}
       s_1 &        0 &        0 &        0 &        0 &        0 \\
\ell_{21} &       s_2 &        0 &        0 &        0 &        0 \\
\ell_{31} & \ell_{32} &       s_3 &        0 &        0 &        0 \\
\ell_{41} & \ell_{42} & \ell_{43} &       s_4 &        0 &        0 \\
\ell_{51} & \ell_{52} & \ell_{53} & \ell_{54} &       s_5 &        0 \\
\ell_{61} & \ell_{62} & \ell_{63} & \ell_{64} & \ell_{65} &       s_6
\end{array}\right),
\end{equation}
where $s_1, s_2, s_3$ are spatial scales, $s_4, s_5, s_6$ are query dimension scales, and $\ell_{ij} \in [-1,1]$ are bounded correlations.

% \resizebox{\textwidth}{!}{%
% \begin{equation}    
% \boldsymbol{\Sigma} = \left(\begin{array}{rrrrrr}
% s_1^2 
% &  
% &  
% &  
% &  
% &  \\
% s_1 \ell_{21} 
% & \ell_{21}^2 + s_2^2 
% &  
% &  
% &  
% &  \\
% s_1 \ell_{31} 
% & \ell_{21}\ell_{31} + s_2\ell_{32} 
% & \ell_{31}^2 + \ell_{32}^2 + s_3^2 
% &  
% &  
% &  \\
% s_1 \ell_{41} 
% & \ell_{21}\ell_{41} + s_2\ell_{42} 
% & \ell_{31}\ell_{41} + \ell_{32}\ell_{42} + s_3\ell_{43} 
% & \ell_{41}^2 + \ell_{42}^2 + \ell_{43}^2 + s_4^2 
% &  
% &  \\
% s_1 \ell_{51} 
% & \ell_{21}\ell_{51} + s_2\ell_{52} 
% & \ell_{31}\ell_{51} + \ell_{32}\ell_{52} + s_3\ell_{53} 
% & \ell_{41}\ell_{51} + \ell_{42}\ell_{52} + \ell_{43}\ell_{53} + s_4\ell_{54} 
% & \ell_{51}^2 + \ell_{52}^2 + \ell_{53}^2 + \ell_{54}^2 + s_5^2 
% &  \\
% s_1 \ell_{61} 
% & \ell_{21}\ell_{61} + s_2\ell_{62} 
% & \ell_{31}\ell_{61} + \ell_{32}\ell_{62} + s_3\ell_{63} 
% & \ell_{41}\ell_{61} + \ell_{42}\ell_{62} + \ell_{43}\ell_{63} + s_4\ell_{64} 
% & \ell_{51}\ell_{61} + \ell_{52}\ell_{62} + \ell_{53}\ell_{63} + \ell_{54}\ell_{64} + s_5\ell_{65} 
% & \ell_{61}^2 + \ell_{62}^2 + \ell_{63}^2 + \ell_{64}^2 + \ell_{65}^2 + s_6^2
% \end{array}\right)
% \end{equation}}

The displacement block $\boldsymbol{\Sigma}_{pq} = (\mathbf{L}\mathbf{L}^\top)_{1:3,4:6}$ emerges from matrix multiplication. Computing via row-column dot products yields:
\begin{equation}
\boldsymbol{\Sigma}_{pq} = \left(\begin{array}{rrr}
                           s_1 \ell_{41} &                            s_1 \ell_{51} &                            s_1 \ell_{61} \\
     \ell_{21}\ell_{41} + s_2 \ell_{42} &      \ell_{21}\ell_{51} + s_2 \ell_{52} &      \ell_{21}\ell_{61} + s_2 \ell_{62} \\
\ell_{31}\ell_{41} + \ell_{32}\ell_{42} + s_3 \ell_{43} & \ell_{31}\ell_{51} + \ell_{32}\ell_{52} + s_3 \ell_{53} & \ell_{31}\ell_{61} + \ell_{32}\ell_{62} + s_3 \ell_{63}
\end{array}\right).
\end{equation}
Note that $\boldsymbol{\Sigma}_{pq}$ contains only spatial scales $s_1, s_2, s_3$ and bounded correlations $\ell_{ij}$ -- the query dimension scales $s_4, s_5, s_6$ do \textit{not} appear. This occurs because displacement describes how spatial position changes with query deviation, which should scale with the Gaussian's \textit{spatial} extent, not its query dimension spread.

In general, each element of $\boldsymbol{\Sigma}_{pq}$ can be expressed as:
\begin{equation}
[\boldsymbol{\Sigma}_{pq}]_{ij} = \sum_{k=1}^{i} L_{ik} L_{(3+j),k} = \sum_{k=1}^{i-1} \ell_{ik} \ell_{(3+j),k} + s_i \cdot \ell_{(3+j),i}.
\end{equation}

Cross-terms from spatial correlations $\ell_{ik}$ (where $k < i$) accumulate in the displacement. The displacement in row $i$ satisfies:
\begin{equation}
\|\boldsymbol{\Sigma}_{pq}^{(i)}\|_2 \leq \sqrt{C}(i-1 + s_i).
\end{equation}
For large $s_i$, this approximates $s_i \sqrt{C}$, but the constant offset $(i-1)\sqrt{C}$ can dominate for small Gaussians. This bound also introduces an asymmetry across spatial dimensions: even when $s_1 = s_2 = s_3$, row $i$ accumulates more cross-terms than row 1, allowing arbitrarily different displacement magnitudes depending on the spatial dimension ordering. While N-DGS uses zero initialization of $\ell_{ij}$ to ensure displacement initially scales with $s_i$, training can learn non-zero correlations that introduce both the offset and this asymmetry.

\textbf{Our method: Explicit coupling via scalar scaling.} We directly parameterize $\mathbf{V}_{pq} = \bar{s} \boldsymbol{\Theta}_{pq}$, where $\bar{s} = \text{mean}(s_x, s_y, s_z)$ is the average spatial scale and $\boldsymbol{\Theta}_{pq} \in \mathbb{R}^{3 \times C}$ are learnable parameters normalized to unit Frobenius norm ($\|\boldsymbol{\Theta}_{pq}\|_F = 1$). This gives:
\begin{equation}
\|\mathbf{V}_{pq}^{(i)}\|_2 = \bar{s} \|\boldsymbol{\Theta}_{i,:}\|_2 \leq \bar{s}.
\end{equation}
Each element in row $i$ scales uniformly with $\bar{s}$ with no cross-terms from other spatial dimensions, ensuring displacement scales with the average Gaussian size even for tiny Gaussians. The unit Frobenius norm constraint ensures bounded displacements while allowing flexible direction learning.

\textbf{Key difference.} Both methods bound displacement by spatial scale, but differ fundamentally:
\begin{itemize}[leftmargin=*]
\item \textbf{N-DGS}: Displacement emerges \textit{implicitly} from Cholesky structure with bound $\sqrt{C}(i-1 + s_i)$. Cross-terms can allow disproportionately large displacements for small Gaussians and introduce asymmetry across spatial dimensions. For highly anisotropic Gaussians (e.g., $s_x \gg s_y, s_z$), displacement along each axis $i$ is constrained by its individual scale $s_i$, allowing excessive motion along large axes while limiting movement along smaller dimensions.
\item \textbf{\modelname{}}: Displacement is \textit{explicitly} parameterized with tight bound $\bar{s}$, ensuring proper scaling across all Gaussian sizes with uniform treatment of all spatial dimensions. Critically, using the mean $\bar{s}$ enables balanced displacement control: it allows sufficient movement along smaller dimensions while preventing excessive motion along larger axes, ensuring the Gaussian can move uniformly in all directions regardless of anisotropy.
\end{itemize}

Empirically, removing spatial scaling (ablation: \textit{d}GS w/o $\bar{s}$) degrades PSNR by 2.59 dB on NeRF Synthetic (Table~\ref{tab:ablation}), confirming that displacement scaling is critical. Our direct parameterization provides tighter bounds, avoids dimension-ordering artifacts, and enables learnable per-dimension coupling strength control via $\Lambda$.


\section{Detailed Per-Scene Results}
\label{sec:detailed_results}

This section provides comprehensive per-scene breakdowns for all experiments. Table~\ref{tab:detailed_combined} presents quality metrics (PSNR, SSIM, LPIPS) for each scene across all five datasets: three static (NeRF Synthetic~\cite{mildenhall2021nerf}, Mip-NeRF 360~\cite{barron2021mipnerf}, 6DGS-PBR~\cite{gao20246dgs}) and two dynamic (D-NeRF~\cite{pumarola2020dnerfneuralradiancefields}, 7DGS-PBR~\cite{gao20257dgs}), while Table~\ref{tab:efficiency_combined} reports efficiency metrics (Gaussian count, training time, rendering FPS). These detailed results demonstrate the consistency of our method's performance across diverse scene types and complexity levels.


% Combined Detailed Quantitative Comparison Table
\begin{table*}[bp]
  \centering
  \caption{\small{Detailed quantitative comparison across all datasets at 30K iterations. N-DGS, dGS-O, and dGS are compared on static benchmarks (NeRF Synthetic, Mip-NeRF 360, 6DGS-PBR) and dynamic benchmarks (7DGS-PBR, D-NeRF).}} %Note, NeRF Synthetic and 7DGS-PBR use decoupled opacity and position.
  \resizebox{1.0\linewidth}{!}{
        \begin{tabular}{@{\hskip 2pt}cc|@{\hskip 2pt}l@{\hskip 2pt}|@{\hskip 2pt}r@{\hskip 2pt}r@{\hskip 2pt}r|@{\hskip 2pt}r@{\hskip 2pt}r@{\hskip 2pt}r|@{\hskip 2pt}r@{\hskip 2pt}r@{\hskip 2pt}r@{\hskip 2pt}}

    \toprule

    \multicolumn{2}{@{\hskip 2pt}c|@{\hskip 2pt}}{\multirow{2}[4]{*}{Dataset}} & \multicolumn{1}{c|@{\hskip 2pt}}{\multirow{2}[4]{*}{Scene}}
        & \multicolumn{3}{c|@{\hskip 2pt}}{N-DGS}
        & \multicolumn{3}{c|@{\hskip 2pt}}{\textit{d}GS-O (\textbf{Ours})}
        & \multicolumn{3}{@{\hskip 2pt}c}{\textit{d}GS (\textbf{Ours})} \\
\cmidrule{4-12}
          &&       & \multicolumn{1}{@{\hskip 2pt}c@{\hskip 2pt}}{PSNR$\uparrow$}
                  & \multicolumn{1}{@{\hskip 2pt}c@{\hskip 2pt}}{SSIM$\uparrow$}
                  & \multicolumn{1}{c|@{\hskip 2pt}}{LPIPS$\downarrow$}
                  & \multicolumn{1}{@{\hskip 2pt}c@{\hskip 2pt}}{PSNR$\uparrow$}
                  & \multicolumn{1}{@{\hskip 2pt}c@{\hskip 2pt}}{SSIM$\uparrow$}
                  & \multicolumn{1}{c|@{\hskip 2pt}}{LPIPS$\downarrow$}
                  & \multicolumn{1}{@{\hskip 2pt}c@{\hskip 2pt}}{PSNR$\uparrow$}
                  & \multicolumn{1}{@{\hskip 2pt}c@{\hskip 2pt}}{SSIM$\uparrow$}
                  & \multicolumn{1}{@{\hskip 2pt}c@{\hskip 2pt}}{LPIPS$\downarrow$} \\
    \midrule

\multicolumn{1}{@{\hskip 2pt}c}{\multirow{30}[4]{*}{\begin{sideways}STATIC SCENES\end{sideways}}}

    % ---------------- nerf_synthetic ----------------
& \multicolumn{1}{c|@{\hskip 2pt}}{\multirow{9}[4]{*}{\begin{sideways}NeRF Synthetic\end{sideways}}}
    & \texttt{chair}     & 35.75 & 0.986 & 0.013 & \bf 36.36 & \bf 0.988 & \bf 0.012 &  36.26 & \bf 0.988 & \bf 0.012 \\
    & & \texttt{drums}     & 26.75 & 0.956 & 0.036 & 26.80 & 0.957 & 0.035 & \bf 26.92 & \bf 0.958 & \bf 0.034 \\
    & & \texttt{ficus}     & 33.54 & 0.984 & 0.014 & \bf 35.30 & \bf 0.988 & \bf 0.011 & 35.03 & 0.988 & \bf 0.011 \\
    & & \texttt{hotdog}    & 38.10 & 0.985 & 0.022 & \bf 38.13 & \bf 0.986 & \bf 0.020 & \bf 38.15 & \bf 0.986 & \bf 0.020 \\
    & & \texttt{lego}      & 35.63 & 0.981 & 0.019 & \bf 35.87 & \bf 0.982 & \bf 0.017 & 35.79 & \bf 0.982 & \bf 0.017 \\
    & & \texttt{materials} & 30.78 & 0.967 & 0.033 & 31.00 & 0.967 & 0.030 & \bf 31.76 & \bf 0.972 & \bf 0.027 \\
    & & \texttt{mic}       & 35.63 & 0.992 & 0.006 & 36.12 & \bf 0.993 & \bf 0.005 & \bf 36.25 & \bf 0.993 & \bf 0.005 \\
    & & \texttt{ship}      & 30.85 & 0.906 & 0.112 & \bf 31.16 & \bf 0.908 & 0.107 & 31.13 & \bf 0.908 & \bf 0.106 \\
\cmidrule{3-12}
    & & \texttt{avg}
      & 33.38
      & 0.970
      & 0.032
      & 33.84
      & 0.971
      & 0.030
      & \bf 33.91
      & \bf 0.972
      & \bf 0.029 \\
    \cmidrule{2-12}

    % ---------------- mipnerf360 ----------------
& \multicolumn{1}{c|@{\hskip 2pt}}{\multirow{10}[4]{*}{\begin{sideways}Mip-NeRF 360\end{sideways}}}
        & \texttt{bicycle}   & 22.33 & 0.554 & 0.447 & 23.49 & 0.644 & 0.327 & \bf 23.51 & \bf 0.648 & \bf 0.318 \\
        & & \texttt{bonsai}    & 33.20 & 0.940 & 0.208 & 33.22 & 0.949 & 0.180 & \bf 33.46 & \bf 0.951 & \bf 0.175 \\
        & & \texttt{counter}   & 29.26 & 0.895 & 0.236 & 29.73 & 0.917 &  0.185 & \bf 29.93 & \bf 0.920 & \bf 0.181 \\
        & & \texttt{flowers}   & 18.90 & 0.409 & 0.507 &  19.61 & \bf 0.453 & \bf 0.443 & \bf 19.71 & 0.452 & 0.455 \\
        & & \texttt{garden}    & 26.19 & 0.809 & 0.205 & 26.77 & 0.842 &  0.145 & \bf 27.13 & \bf 0.849 & \bf 0.137 \\
        & & \texttt{kitchen}   & 31.17 & 0.922 & 0.137 & 31.72 & 0.930 &  0.121 & \bf 32.04 & \bf 0.933 & \bf 0.116 \\
        & & \texttt{room}      & 31.60 & 0.910 & 0.243 & 31.83 & 0.927 & 0.199 & \bf 32.25 & \bf 0.929 & \bf 0.195 \\
        & & \texttt{stump}     & 23.32 & 0.638 & 0.413 & \bf 26.01 &  0.750 & \bf 0.275 & 26.00 & \bf 0.742 & 0.295 \\
        & & \texttt{treehill}  & 21.66 & 0.571 & 0.448 & 22.21 & 0.603 & 0.393 & \bf 22.31 & \bf 0.606 & \bf 0.389 \\
    \cmidrule{3-12}
        & & \texttt{avg}
          & 26.40
          & 0.739
          & 0.316
          & 27.18
          & 0.779
          & 0.252
          & \bf 27.37
          & \bf 0.781
          & \bf 0.251 \\
\cmidrule{2-12}

  % ---------------- tandt_pbr ----------------
& \multicolumn{1}{c|@{\hskip 2pt}}{\multirow{8}[4]{*}{\begin{sideways}6DGS-PBR\end{sideways}}}
    & \texttt{bunny} & 40.87 & 0.993 & 0.045 & 43.44 & 0.993 & 0.042 & \bf 44.14 & \bf 0.993 & \bf 0.040 \\
    & & \texttt{cloud}     & 41.93 & 0.991 & 0.052 & 43.47 & 0.992 & 0.048 & \bf 43.53 & \bf 0.992 & \bf 0.045 \\
    & & \texttt{dragon}    &  35.69 & 0.947 & 0.039 & 34.42 & 0.914 & 0.051 &  \bf 36.68 &  \bf 0.947 & \bf 0.039 \\
    & & \texttt{explosion} & \bf 40.97 & \bf 0.989 & \bf 0.037 & 37.81 & 0.983 & 0.041 & 38.95 & 0.985 & 0.039 \\
    & & \texttt{ct-scan}     & 33.48 & 0.964 & 0.057 & 31.81 & 0.956 & 0.067 & \bf 34.62 & \bf 0.970 & \bf 0.048 \\
    & & \texttt{smoke}     & \bf 40.02 & \bf 0.990 & \bf 0.041 & 38.20 & 0.989 & 0.049 & 39.28 & \bf 0.990 & 0.047 \\
    & & \texttt{suzanne}   & 26.81 & 0.929 & \bf 0.106 & 26.65 & 0.927 & 0.112 & \bf 27.03 & \bf 0.930 & 0.107 \\
\cmidrule{3-12}
    & & \texttt{avg}
      & 37.11
      & 0.971
      & 0.054
      & 36.54
      & 0.965
      & 0.059
      & \bf 37.75
      & \bf 0.972
      & \bf 0.052 \\

    \midrule
    \midrule

\multicolumn{1}{@{\hskip 2pt}c}{\multirow{16}[4]{*}{\begin{sideways}DYNAMIC SCENES\end{sideways}}}

% & \multirow{7}{*}{\begin{sideways}7DGS-PBR\end{sideways}}
%     & \texttt{cloud}    & \bf 30.31 & \bf 0.956 & \bf 0.072 & 29.34 & 0.955 & 0.075 & 29.60 & 0.955 & 0.074 \\
%     & & \texttt{dust}     & \bf 39.30 & \bf 0.959 & \bf 0.031 & 36.87 & 0.955 & 0.037 & 36.63 & 0.954 & 0.037 \\
%     & & \texttt{flame}    & \bf 33.90 & \bf 0.945 & \bf 0.054 & 30.38 & 0.933 & 0.066 & 30.32 & 0.932 & 0.066 \\
%     & & \texttt{heart}   & \bf 36.09 & \bf 0.987 & \bf 0.018 & 35.14 & 0.985 & 0.021 & 35.25 & 0.985 & 0.020 \\
%     & & \texttt{heart\_1600}   & \bf 32.15 & \bf 0.966 & \bf 0.048 & 31.29 & 0.962 & 0.052 & 31.36 & 0.962 & 0.051 \\
%     & & \texttt{suzanne}  & \bf 27.53 & \bf 0.947 & \bf 0.061 & 27.20 & 0.940 & 0.069 & 27.19 & 0.940 & 0.072 \\
%     \cmidrule{3-12}
%         & & \texttt{avg}
%           & \bf 33.21
%           & \bf 0.960
%           & \bf 0.047
%           & 31.70
%           & 0.955
%           & 0.053
%           & 31.72
%           & 0.955
%           & 0.053 \\
%     \cmidrule{2-12}


& \multirow{7}{*}{\begin{sideways}7DGS-PBR\end{sideways}}
    & \texttt{cloud}    & 29.28 & 0.955 & 0.075 & 29.34 & 0.955 & 0.075 & \bf 29.82 & \bf 0.955 & \bf 0.074 \\
    & & \texttt{dust}     & 36.85 & 0.955 & 0.038 & 36.87 & 0.955 &  0.037 & \bf 38.39 & \bf 0.958 & \bf 0.031 \\
    & & \texttt{flame}    & \bf 31.64 & \bf 0.937 & \bf 0.062 & 30.38 & 0.933 & 0.066 & 31.52 & 0.936 & 0.060 \\
    & & \texttt{heart1}   & 34.66 & 0.983 & 0.023 & 35.14 & \bf 0.985 & 0.021 &  \bf 35.25 & \bf 0.986 & \bf 0.018 \\
    & & \texttt{heart2}   & 30.99 & 0.959 & 0.057 & \bf 31.29 & \bf 0.962 & 0.052 & 31.01 & \bf 0.962 & \bf 0.049 \\
    & & \texttt{suzanne}  & 27.09 & \bf 0.941 & 0.072 & \bf 27.20 & 0.940 & \bf 0.069 & 27.07 & 0.940 & 0.070 \\
    \cmidrule{3-12}
        & & \texttt{avg}
          & 31.75
          & 0.955
          & 0.055
          & 31.70
          & 0.955
          & 0.053
          & \bf 32.18
          & \bf 0.956
          & \bf 0.050 \\
    \cmidrule{2-12}

& \multirow{9}{*}{\begin{sideways}D-NeRF\end{sideways}}
        & \texttt{b.balls}  & \bf 34.00 & \bf 0.981 & 0.025 & 32.52 & 0.978 & 0.028 &  33.59 & \bf 0.981 & \bf 0.023 \\
        & & \texttt{h.warrior}& 32.76 & 0.932 & 0.093 & 33.54 & 0.941 & 0.085 & \bf 34.64 & \bf 0.947 & \bf 0.073 \\
        & & \texttt{hook}     & 31.15 & 0.960 & 0.044 & 28.55 & 0.937 & 0.062 & \bf 32.94 & \bf 0.970 & \bf 0.030 \\
        & & \texttt{j.jacks}  & \bf 31.22 & \bf 0.967 & 0.043 & 28.62 & 0.952 & 0.054 & 30.86 & 0.965 & \bf 0.039 \\
        & & \texttt{lego}     & \bf 28.33 & \bf 0.943 & 0.055 & 27.13 & 0.939 &  0.055 & 28.06 & 0.942 & \bf 0.054 \\
        & & \texttt{mutant}   & 39.38 & 0.993 & 0.007 & 34.55 & 0.981 & 0.019 & \bf 39.84 & \bf 0.993 & \bf 0.006 \\
        & & \texttt{standup}  & 38.50 & 0.988 & 0.013 & 33.02 & 0.972 & 0.029 & \bf 38.76 & \bf 0.990 & \bf 0.009 \\
        & & \texttt{trex}     & 29.68 & 0.978 & 0.023 & 29.42 & 0.975 & 0.022 & \bf 30.08 & \bf 0.978 & \bf 0.021 \\
    \cmidrule{3-12}
        & & \texttt{avg}
          & 33.13
          & 0.968
          & 0.038
          & 30.92
          & 0.959
          & 0.044
          & \bf 33.59
          & \bf 0.971
          & \bf 0.032 \\
    \bottomrule
    \end{tabular}
    }
  \label{tab:detailed_combined}
\end{table*}


% Combined Detailed Efficiency Comparison Table
\begin{table*}[tbp]
  \centering
  \caption{\small{Detailed efficiency comparison across all datasets at 30K iterations. N-DGS, \textit{d}GS-O, and \textit{d}GS are compared showing number of Gaussians (K), training time (minutes), and rendering FPS on static benchmarks (NeRF Synthetic, Mip-NeRF 360, 6DGS-PBR) and dynamic benchmarks (7DGS-PBR, D-NeRF).}} %Note, NeRF Synthetic and 7DGS-PBR use decoupled opacity and position.
  \resizebox{0.85\linewidth}{!}{
        \begin{tabular}{@{\hskip 2pt}cc|@{\hskip 2pt}l@{\hskip 2pt}|@{\hskip 2pt}r@{\hskip 2pt}r@{\hskip 2pt}r|@{\hskip 2pt}r@{\hskip 2pt}r@{\hskip 2pt}r|@{\hskip 2pt}r@{\hskip 2pt}r@{\hskip 2pt}r@{\hskip 2pt}}

    \toprule

    \multicolumn{2}{@{\hskip 2pt}c|@{\hskip 2pt}}{\multirow{2}[4]{*}{Dataset}} & \multicolumn{1}{c|@{\hskip 2pt}}{\multirow{2}[4]{*}{Scene}}
        & \multicolumn{3}{c|@{\hskip 2pt}}{N-DGS}
        & \multicolumn{3}{c|@{\hskip 2pt}}{\textit{d}GS-O (\textbf{Ours})}
        & \multicolumn{3}{@{\hskip 2pt}c}{\textit{d}GS (\textbf{Ours})} \\
\cmidrule{4-12}
          &&       & \multicolumn{1}{@{\hskip 2pt}c@{\hskip 2pt}}{\#Gauss$\downarrow$}
                  & \multicolumn{1}{@{\hskip 2pt}c@{\hskip 2pt}}{Time$\downarrow$}
                  & \multicolumn{1}{c|@{\hskip 2pt}}{FPS$_{\text{CUDA}}$$\uparrow$}
                  & \multicolumn{1}{@{\hskip 2pt}c@{\hskip 2pt}}{\#Gauss$\downarrow$}
                  & \multicolumn{1}{@{\hskip 2pt}c@{\hskip 2pt}}{Time$\downarrow$}
                  & \multicolumn{1}{c|@{\hskip 2pt}}{FPS$_{\text{CUDA}}$$\uparrow$}
                  & \multicolumn{1}{@{\hskip 2pt}c@{\hskip 2pt}}{\#Gauss$\downarrow$}
                  & \multicolumn{1}{@{\hskip 2pt}c@{\hskip 2pt}}{Time$\downarrow$}
                  & \multicolumn{1}{@{\hskip 2pt}c@{\hskip 2pt}}{FPS$_{\text{CUDA}}$$\uparrow$} \\
    \midrule

\multicolumn{1}{@{\hskip 2pt}c}{\multirow{30}[4]{*}{\begin{sideways}STATIC SCENES\end{sideways}}}

    % ---------------- nerf_synthetic ----------------
& \multicolumn{1}{c|@{\hskip 2pt}}{\multirow{9}[4]{*}{\begin{sideways}NeRF Synthetic\end{sideways}}}
    & \texttt{chair}     & \bf 178 & 6.11 & 635.27 & 206 & \bf 5.54 & \bf 713.54 & 199 & 5.86 & 677.38 \\
    & & \texttt{drums}     & \bf 193 & 6.19 & 476.51 & 239 & \bf 5.62 & \bf 683.86 & 242 & 6.04 & 646.90 \\
    & & \texttt{ficus}     & \bf 117 & 4.58 & 840.56 & 147 & \bf 4.27 & \bf 973.46 & 136 & 4.53 & 900.59 \\
    & & \texttt{hotdog}    & 102 & 6.04 & 871.45 & \bf 98 & \bf 5.56 & \bf 982.05 & 104 & 6.02 & 876.19 \\
    & & \texttt{lego}      & \bf 186 & 7.06 & 480.97 & 220 & \bf 5.67 & \bf 697.44 & 225 & 6.23 & 669.15 \\
    & & \texttt{materials} & \bf 136 & 5.38 & 767.74 & 182 & \bf 5.09 & 759.30 & 180 & 5.54 & 689.29 \\
    & & \texttt{mic}       & \bf 188 & 5.84 & 486.29 & 260 & 6.21 & \bf 625.05 & 260 & 11.17 & 594.08 \\
    & & \texttt{ship}      & 271 & 9.60 & 365.77 & \bf 256 & \bf 8.68 & \bf 605.15 & 259 & 9.28 & 585.17 \\
\cmidrule{3-12}
    & & \texttt{avg}
      & \bf 171
      & 6.35
      & 615.57
      & 201
      & \bf 5.83
      & \bf 754.98
      & 201
      & 6.84
      & 704.84 \\
    \cmidrule{2-12}

    % ---------------- mipnerf360 ----------------
& \multicolumn{1}{c|@{\hskip 2pt}}{\multirow{10}[4]{*}{\begin{sideways}Mip-NeRF 360\end{sideways}}}
        & \texttt{bicycle}   & \bf 2243 & \bf 32.54 & 70.31 & 4452 & 33.46 & \bf 144.01 & 4644 & 38.30 & 118.20 \\
        & & \texttt{bonsai}    & \bf 751 & 23.49 & 186.80 & 1048 & \bf 20.02 & \bf 419.45 & 1091 & 21.77 & 359.76 \\
        & & \texttt{counter}   & \bf 555 & 22.79 & 237.17 & 880 & \bf 22.41 & \bf 430.22 & 902 & 23.98 & 380.55 \\
        & & \texttt{flowers}   & \bf 1547 & 26.28 & 100.84 & 2587 & \bf 24.43 & \bf 237.63 & 2679 & 26.27 & 196.76 \\
        & & \texttt{garden}    & \bf 2674 & 38.90 & 58.64 & 3552 & \bf 29.36 & \bf 174.37 & 3767 & 33.70 & 141.57 \\
        & & \texttt{kitchen}   & \bf 770 & 27.28 & 178.89 & 1148 & \bf 24.95 & \bf 354.03 & 1175 & 26.41 & 312.45 \\
        & & \texttt{room}      & \bf 792 & 24.62 & 177.95 & 1286 & \bf 23.81 & \bf 351.34 & 1356 & 25.75 & 298.32 \\
        & & \texttt{stump}     & \bf 644 & 21.29 & 212.22 & 3272 & \bf 27.56 & 199.15 & 3343 & 30.01 & 165.52 \\
        & & \texttt{treehill}  & \bf 1481 & 26.62 & 104.35 & 3042 & \bf 24.22 & \bf 203.51 & 3209 & 27.64 & 167.23 \\
    \cmidrule{3-12}
        & & \texttt{avg}
          & \bf 1273
          & 27.09
          & 147.46
          & 2363
          & \bf 25.58
          & \bf 279.30
          & 2463
          & 28.21
          & 237.82 \\
\cmidrule{2-12}

  % ---------------- tandt_pbr ----------------
& \multicolumn{1}{c|@{\hskip 2pt}}{\multirow{8}[4]{*}{\begin{sideways}6DGS-PBR\end{sideways}}}
    & \texttt{bunny} & \bf 24 & 33.69 & \bf 718.47 & 63 & 42.62 & 689.88 & 62 & 45.03 & 642.86 \\
    & & \texttt{cloud}     & \bf 17 & 32.61 & 597.10 & 62 & 39.08 & \bf 700.35 & 60 & 41.31 & 654.93 \\
    & & \texttt{dragon}    & \bf 93 & 12.76 & \bf 855.89 & 118 & 12.87 & 787.64 & 134 & 13.28 & 753.36 \\
    & & \texttt{explosion} & \bf 35 & 22.13 & \bf 875.48 & 72 & 29.54 & 813.28 & 73 & 40.77 & 753.18 \\
    & & \texttt{ct-scan}     & \bf 171 & 6.57 & 689.47 & 221 & \bf 6.44 & \bf 796.54 & 299 & 7.45 & 695.69 \\
    & & \texttt{smoke}     & \bf 52 & 46.48 & \bf 608.02 & 72 & 55.67 & 568.30 & 74 & 58.15 & 547.93 \\
    & & \texttt{suzanne}   & \bf 201 & 33.16 & 404.34 & 221 & \bf 28.80 & \bf 518.17 & 233 & 30.58 & 526.10 \\
\cmidrule{3-12}
    & & \texttt{avg}
      & \bf 85
      & 26.77
      & 678.40
      & 118
      & 30.72
      & \bf 696.31
      & 134
      & 33.80
      & 653.44 \\

    \midrule
    \midrule

\multicolumn{1}{@{\hskip 2pt}c}{\multirow{16}[4]{*}{\begin{sideways}DYNAMIC SCENES\end{sideways}}}

& \multirow{7}{*}{\begin{sideways}7DGS-PBR\end{sideways}}
    & \texttt{cloud}    & \bf 48 & 76.87 & 950.04 & 63 & \bf 75.47 & \bf 953.38 & 70 & 76.11 & 892.36 \\
    & & \texttt{dust}     & \bf 16 & 13.11 & 1052.85 & 23 & \bf 12.06 & \bf 1165.85 & 30 & 16.89 & 990.11 \\
    & & \texttt{flame}    & \bf 22 & \bf 14.91 & \bf 1063.87 & 23 & 16.87 & 1015.92 & 40 & 19.30 & 926.17 \\
    & & \texttt{heart1}   & \bf 140 & 30.75 & 593.96 & 182 & \bf 28.75 & \bf 697.33 & 188 & 30.08 & 639.83 \\
    & & \texttt{heart2}   & \bf 181 & 22.75 & 432.58 & 237 & \bf 20.12 & \bf 729.08 & 261 & 22.60 & 662.23 \\
    & & \texttt{suzanne}  & \bf 564 & 70.47 & 165.70 & 837 & \bf 54.23 & \bf 449.89 & 767 & 56.75 & 411.89 \\
    \cmidrule{3-12}
        & & \texttt{avg}
          & \bf 162
          & 38.14
          & 709.83
          & 227
          & \bf 34.58
          & \bf 835.24
          & 226
          & 36.96
          & 753.77 \\
    \cmidrule{2-12}

& \multirow{9}{*}{\begin{sideways}D-NeRF\end{sideways}}
        & \texttt{b.balls}  & 179 & 22.73 & 398.65 & \bf 103 & \bf 16.09 & \bf 770.99 & 128 & 17.72 & 636.37 \\
        & & \texttt{h.warrior}& \bf 13 & 8.25 & \bf 1151.03 & 17 & \bf 7.87 & 1141.34 & 18 & 9.23 & 1068.65 \\
        & & \texttt{hook}     & \bf 37 & \bf 9.71 & \bf 1042.00 & 55 & 11.42 & 1027.26 & 75 & 14.58 & 865.61 \\
        & & \texttt{j.jacks}  & \bf 25 & \bf 9.61 & 1069.25 & 31 & 10.67 & \bf 1101.29 & 71 & 14.48 & 899.56 \\
        & & \texttt{lego}     & \bf 114 & 21.09 & 583.20 & 143 & \bf 17.93 & \bf 700.47 & 125 & 19.28 & 685.07 \\
        & & \texttt{mutant}   & \bf 63 & 14.74 & 829.30 & 67 & \bf 13.27 & \bf 930.20 & 79 & 15.47 & 810.36 \\
        & & \texttt{standup}  & \bf 21 & 9.71 & 1074.73 & 26 & \bf 9.28 & \bf 1140.56 & 51 & 14.00 & 867.39 \\
        & & \texttt{trex}     & \bf 101 & \bf 16.77 & 667.96 & 118 & 19.09 & \bf 778.03 & 134 & 19.69 & 714.27 \\
    \cmidrule{3-12}
        & & \texttt{avg}
          & \bf 69
          & 14.08
          & 852.01
          & 70
          & \bf 13.20
          & \bf 948.77
          & 85
          & 15.56
          & 818.41 \\
    \bottomrule
    \end{tabular}
    }
  \label{tab:efficiency_combined}
\end{table*}


\section{Combining \modelname{} with Alternative Kernels}
\label{sec:dbs}

\modelname{} targets the \textit{conditioning parameterization} of N-dimensional Gaussians, while UBS~\citep{liu2025universal} proposes an alternative \textit{kernel shape} (Beta kernels) for the same N-dimensional framework. These contributions are orthogonal: \modelname{}'s direct parameterization can replace UBS's covariance-derived conditioning while retaining UBS's Beta kernel. We first compare both methods directly, then present a concrete combined formulation.

\subsection{Direct Comparison: \modelname{} vs UBS}

Table~\ref{tab:dgs_vs_ubs} compares \modelname{} and UBS on shared static benchmarks using MCMC densification. UBS achieves higher quality on NeRF Synthetic (+0.58 dB) and Mip-NeRF 360 (+0.99 dB), while \modelname{} is comparable on 6DGS-PBR (+0.24 dB). Importantly, \modelname{} is consistently \textbf{$\sim$2$\times$ faster} in rendering across all datasets. This motivates combining both: UBS's Beta kernel quality with \modelname{}'s efficient conditioning.

\begin{table}[h]
\centering
\caption{Direct comparison of \modelname{} and UBS on shared static benchmarks (MCMC densification). Both use the same Gaussian budgets.}
\label{tab:dgs_vs_ubs}
\setlength{\tabcolsep}{4pt}
\resizebox{0.5\textwidth}{!}{%
\begin{tabular}{@{}cl|cccc@{}}
\toprule
Dataset & Method & PSNR $\uparrow$ & SSIM $\uparrow$ & LPIPS $\downarrow$ & FPS $\uparrow$ \\
\midrule
\multirow{2}{*}{\makecell{\texttt{NeRF}\\\texttt{Synthetic}}}
& \textit{d}GS & 34.34 & 0.973 & 0.026 & \textbf{594} \\
& UBS & \textbf{34.92} & \textbf{0.975} & \textbf{0.026} & 346 \\
\midrule
\multirow{2}{*}{\makecell{\texttt{Mip-NeRF}\\\texttt{360}}}
& \textit{d}GS & 27.67 & 0.794 & 0.227 & \textbf{201} \\
& UBS & \textbf{28.66} & \textbf{0.840} & \textbf{0.184} & 100 \\
\midrule
\multirow{2}{*}{\texttt{6DGS-PBR}}
& \textit{d}GS & \textbf{40.34} & 0.976 & \textbf{0.038} & \textbf{551} \\
& UBS & 40.10 & \textbf{0.979} & 0.039 & 290 \\
\bottomrule
\end{tabular}}
\end{table}

\subsection{dBS: Direct Beta Splatting}

We propose \textbf{dBS} (direct Beta Splatting), which combines \modelname{}'s direct conditioning with UBS's Beta kernel. The key idea is to replace UBS's covariance-derived conditioning ($\boldsymbol{\Sigma}_{qq}^{-1}$-based) with \modelname{}'s direct parameterization while keeping the per-dimension Beta bandwidth control.

\textbf{UBS} (current formulation):
\begin{align}
d_i &= \tanh\!\left((\hat{\mathbf{L}}_q^{-1} \boldsymbol{\delta})_i^2 / 2\right), \\
\alpha_{\text{cond}} &= \alpha \prod_{i=1}^{C} (1 - d_i)^{4\beta_i}, \\
\boldsymbol{\mu}_{\text{cond}} &= \boldsymbol{\mu}_p + \boldsymbol{\Sigma}_{pq} \boldsymbol{\Sigma}_{qq}^{-1} \operatorname{diag}(\boldsymbol{\beta}_q) \boldsymbol{\delta}, \\
\boldsymbol{\Sigma}_{\text{cond}} &= \boldsymbol{\Sigma}_{pp} - \boldsymbol{\Sigma}_{pq} \boldsymbol{\Sigma}_{qq}^{-1} \operatorname{diag}(\boldsymbol{\beta}_q) \boldsymbol{\Sigma}_{qp}.
\end{align}

\textbf{dBS} replaces covariance-derived conditioning with \modelname{}'s direct parameterization:
\begin{align}
d_i &= \tanh\!\left((\mathbf{L}^\top \boldsymbol{\delta})_i^2 / 2\right), \\
\alpha_{\text{cond}} &= \alpha \prod_{i=1}^{C} (1 - d_i)^{4\beta_i}, \\
\boldsymbol{\mu}_{\text{cond}} &= \boldsymbol{\mu}_p + \mathbf{V}_{pq} \cdot \operatorname{diag}(\boldsymbol{\beta}_q) \cdot \mathbf{V}_{qq} \cdot \boldsymbol{\delta}, \\
\boldsymbol{\Sigma}_{\text{cond}} &= \boldsymbol{\Sigma}_{pp} \quad \text{(no correction)}.
\end{align}

Key differences from UBS:
\begin{itemize}[leftmargin=*]
\item \textbf{No inversion}: $\mathbf{L}$ is directly parameterized as a Cholesky precision factor, replacing $\hat{\mathbf{L}}_q^{-1}$.
\item \textbf{Direct displacement}: $\mathbf{V}_{pq}$ replaces $\boldsymbol{\Sigma}_{pq}\boldsymbol{\Sigma}_{qq}^{-1}$, eliminating the regression matrix.
\item \textbf{Unified coupling}: Per-dimension $\boldsymbol{\beta}_q$ from UBS replaces \modelname{}'s $\boldsymbol{\Lambda}$, controlling both opacity bandwidth and position coupling in a single parameter per dimension.
\item \textbf{No covariance correction}: As validated in our experiments, direct $\boldsymbol{\Sigma}_{pp}$ suffices.
\end{itemize}

dBS combines UBS's per-dimension Beta bandwidth control (enabling different bandwidths for diffuse vs specular, static vs dynamic) with \modelname{}'s efficient direct conditioning (5--6$\times$ faster slicing). Implementation and evaluation of dBS is a concrete next step for this line of work.


\end{document}
